\documentclass[11pt]{amsart}
\usepackage[a4paper,margin=26mm]{geometry}
\usepackage[T1]{fontenc}
\usepackage{lmodern,mathtools,amssymb,booktabs,microtype}
\usepackage{xcolor}
\usepackage[colorlinks=true,linkcolor=blue!50!black,citecolor=blue!50!black,urlcolor=blue!50!black]{hyperref}
\hypersetup{pdftitle={Contact and radial poles of dilated Euler-product sections},pdfauthor={K. Srinivasa Raghava}}
\newtheorem{theorem}{Theorem}[section]
\newtheorem{proposition}[theorem]{Proposition}
\newtheorem{lemma}[theorem]{Lemma}
\newtheorem{corollary}[theorem]{Corollary}
\theoremstyle{definition}
\newtheorem{example}[theorem]{Example}
\theoremstyle{remark}
\newtheorem{remark}[theorem]{Remark}
\numberwithin{equation}{section}
\newcommand{\Q}{\mathbb Q}\newcommand{\Z}{\mathbb Z}\newcommand{\C}{\mathbb C}
\newcommand{\ord}{\operatorname{ord}}\newcommand{\Res}{\operatorname{Res}}
\newcommand{\sech}{\operatorname{sech}}
\newcommand{\dd}{\,d}
\setlength{\emergencystretch}{2em}
\title[Root-of-Unity Sections of Dilated Euler-Product Quotients]{Root-of-Unity Sections of Dilated Euler-Product Quotients}
\author{K. Srinivasa Raghava}
\address{Pie Mathematics Association}
\email{srinivasaraghavak@gmail.com}
\dedicatory{This work is dedicated to the enduring spirit of Srinivasa Ramanujan.}
\date{September 7, 2026}
\subjclass[2020]{Primary 11F20, 33D15; Secondary 11F11, 11G15, 30B10, 41A60}
\keywords{Euler product, dilate quotient, root-of-unity section, Lambert series, Dedekind eta-function, complex multiplication, Poisson summation}
\begin{document}
\begin{abstract}
We study root-of-unity sections of quotients formed from dilated Euler products. At the boundary point $q=1$, odd sections tend to nonzero cyclotomic constants, whereas even sections have infinitely many simple real poles. We determine the locations, spacing, and residues of all sufficiently late poles. In the cubic case we find the leading oscillatory error in an elementary approximation and prove that its exponential scale is sharp. At the origin we determine exact contact, including a gap-dependent parity law, and integral divisibility with the first nonzero coefficient determined for every modulus $|k-1|\ge2$. For Euler powers $k\ge2$, a moving formal pole controls the section error at small fixed nomes; a real zero-free interval gives a further convergence regime. For the twenty-fourth power we prove explicit CM approximation bounds and a signed asymptotic as the dilation increases.
\end{abstract}
\maketitle

\section{Introduction}
Let $\Phi(q)=\prod_{n\ge1}(1-q^n)$. For integers $a,d\ge2$ and $k\ge1$, consider
\begin{equation}\label{introW}
 W_{a,d}(q,k)=\frac1d\sum_{z^d=1}
 \frac{\sum_{m\ge0}z^mq^{am+1}\Phi(q^{am+1})^k}
      {\sum_{m\ge0}z^mq^{a(m+1)}\Phi(q^{a(m+1)})^k},
 \qquad
 T(q)=q^{1-a}\frac{\Phi(q)^k}{\Phi(q^a)^k}.
\end{equation}
The quotients are meromorphic for $0<|q|<1$. Their behaviour at $q=0$ and $q=1$ is sharply different.

The dilation $a$, section order $d$, and Euler power $k$ play different roles. Our first summary concerns analytic limits; parameters remain fixed except where a limit is specified.

\begin{theorem}\label{main}
Let $a,d\ge2$ and $k\ge1$ be integers.
\begin{enumerate}
\item As $q\to1^-$, every odd section tends to
\[
 \frac{\sin(\pi/a)}{d\sin(\pi(a-1)/(ad))}.
\]
Every even section has simple real poles $q_j\to1$, containing all its real poles sufficiently near $1$, with
\[
 1-q_j\sim\frac{\pi k}{3aj^2},\qquad
 \Res_{q=q_j}W_{a,d}(q,k)\sim\frac{2k\sin(\pi/a)}{3ad\,j^3}.
\]
\item Let $Q_*(k)$ be the unique solution in $(0,1)$ of
$Q\prod_{n\ge1}(1+Q^n)^k=1$. At each fixed real $0<q\le Q_*(k)^{1/a}$, all sections are finite and converge exponentially to $T(q)$ as $d\to\infty$.
\item Put $q_a=e^{-2\pi/\sqrt a}$. For $k=24$, all sections are finite at $q_a$ and converge exponentially to $a^6$ as $d\to\infty$. For each fixed $d\ge2$,
\[
 W_{a,d}(q_a,24)-a^6
 \sim48\pi(-1)^{d-1}23^{d-2}a^{-1/2}e^{-4\pi(d-1)\sqrt a}
 \quad(a\to\infty).
\]
\end{enumerate}
\end{theorem}
\begin{proof}
The radial assertions are Theorems~\ref{radial} and \ref{poles}; the fixed-nome convergence interval is Theorem~\ref{zerofree}. The CM bound and signed asymptotic are Theorem~\ref{CMtheorem} and Proposition~\ref{CMleading}.
\end{proof}

The cubic refinement in Theorem~\ref{cubicpoles} determines an oscillatory leading error whose scaled limit points fill $[-2,2]$. The next summary describes the formal contact and its arithmetic.

\begin{theorem}\label{mainformal}
Let $a,d\ge2$ and $k\ge1$ be integers. Set
\[
 E_0=2a(d-1)+2,\qquad
 E_d=\begin{cases}5a(d-1)/2+2,&d\text{ odd},\\a(5d-6)/2+3,&d\text{ even}.\end{cases}
\]
For $k\ge2$, the defect begins with
\[
 W_{a,d}(q,k)-T(q)=k(1-k)^{d-1}q^{E_0}+O(q^{E_0+1});
\]
for $k=1$ it begins with $(-1)^{d-1}q^{E_d}$. For $k\ge3$, write $h=k-1$. If $0\le r\le E_d-E_0$, then
\[
 h^{\,d-2-\lfloor r/a\rfloor-\lfloor(r-1)/a\rfloor}
 \ \bigm|\ [q^{E_0+r}]\bigl(W_{a,d}(q,k)-T(q)\bigr).
\]
The exponent of $h$ is nonnegative throughout the displayed range. For $h\ge3$, the first nonzero coefficient modulo $h$ occurs at $E_d$ and equals
$(-1)^{d-1}\{-k(k-3)/2\}^{\lfloor d/2\rfloor}$. For $k=3$, the first nonzero coefficient modulo two is one, at order $3a(d-1)+2$ for odd $d$ and $a(3d-4)+4$ for even $d$.
\end{theorem}
\begin{proof}
Apply Corollaries~\ref{etacontact} and \ref{PhiSecond}, Theorem~\ref{adic}, and Corollary~\ref{modtwoendpoint}.
\end{proof}

When $q=e^{2\pi i\tau}$, the points $q=0$ and $q=1$ correspond to the cusps $\infty$ and $0$. We use this cusp terminology for the modular target. For $\Delta(\tau)=q\Phi(q)^{24}$ the target is
\begin{equation}\label{targetmod}
 T(q)=\frac{\Delta(\tau)}{\Delta(a\tau)}\qquad(k=24).
\end{equation}
It is a modular unit on $\Gamma_0(a)$: both factors have weight twelve, have no zeros in the upper half-plane, and are meromorphic at the cusps. In contrast, the even sections are not meromorphic modular functions on any finite-index subgroup with meromorphic behaviour at the cusps. In fact, the eta transformation gives
\begin{equation}\label{targetcusp}
 T(e^{-t/a})\sim a^{k/2}
       \exp\!\left(-\frac{\pi^2k(a-1)}{6t}\right)
       \quad(t\downarrow0),
\end{equation}
whereas the odd-section limit in Theorem~\ref{main} is nonzero. Formula~\eqref{targetcusp} follows directly from Lemma~\ref{Gschwartz}. Thus the high contact at one cusp coexists with a different limiting value, or with accumulating poles, at the other.

For the formal arguments we allow a general series $F(q)=1+\sum_{n\ge1}f_nq^n$ over a field $K$ of characteristic zero, except where coefficient identities are explicitly reduced in positive characteristic. Define
\begin{equation}\label{Rdef}
 R(X,q)=\frac{U(X,q)}{V(X,q)}=\sum_{m\ge0}C_m(q)X^m,
 \qquad
 \begin{aligned}
 U(X,q)&=\sum_{m\ge0}X^mF(q^{am+1}),\\
 V(X,q)&=\sum_{m\ge0}X^mF(q^{a(m+1)}).
 \end{aligned}
\end{equation}
Since $V(0,q)=F(q^a)$ is a unit, $R\in K[[q]][[X]]$. For $d\ge1$, put
\begin{equation}\label{Pdef}
 P_d(q)=q^{1-a}\sum_{j\ge0}C_{jd}(q)q^{ajd},
 \qquad T(q)=q^{1-a}\frac{F(q)}{F(q^a)}.
\end{equation}
This selects exactly the coefficients whose indices are multiples of $d$. In particular $P_1=q^{1-a}R(q^a,q)$, and $P_d\to T$ coefficientwise as $d\to\infty$. The classical multisection filter \cite[Chapter III]{Comtet} gives, after adjoining roots of unity,
\begin{equation}\label{rootfilter}
 P_d(q)=\frac1d\sum_{z^d=1}q^{1-a}R(zq^a,q).
\end{equation}
For $F=\Phi^k$ we write $P_d=W_{a,d}(q,k)$. Polynomial statements in $k$ are interpreted in $\Q[k][[q]]$ through the binomial expansion. At every integer specialization,
\[
 F(q)\in\Z[[q]],\qquad q^{a-1}W_{a,d}(q,k)\in\Z[[q]],\qquad
 W_{a,d}(q,k)\in q^{1-a}\Z[[q]].
\]
Thus unnormalized sections are Laurent series; for example the level-two target starts with $q^{-1}$. For $k=24N$ the normalization gives
\[
 T(q)=q^{(a-1)(N-1)}
 \left(\frac{\Delta(\tau)}{\Delta(a\tau)}\right)^N,
\]
so the additional power of $q$ must be retained.

The present article is a self-contained extension of the construction announced in the author's earlier preprint record \cite{RaghavaEarlier}. Its public abstract treats the case $a=d=2$, including depth-six contact, a modular baseline, and CM approximation. These special cases are retained here as part of the general theory. We develop arbitrary dilation and section order, explicit diagonal and integral divisibility bounds, and the parity-dependent radial pole laws; the cubic case admits a sharper error analysis. The earlier abstract also announces globally dominant pole pairs for several Euler powers. We do not use those assertions: the sampling-variable pole theorem here is local, and the nome-variable analysis determines all sufficiently late real poles. All theorem statements about these sections are proved here, with the classical inputs cited explicitly.

The first nonconstant coefficient of $F$ determines the first defect. Clearing one Lambert denominator then makes each coefficient diagonal rational, with explicit denominator and degree bounds. Rational generating functions are classical \cite{Stanley,FlajoletSedgewick}; the bounds here determine both the higher contact at $f_1=-1$ and a divisibility law modulo powers of $k-1$. The formal pole of these diagonal functions is an actual simple pole for small nonzero $q$, by Theorem~\ref{analyticpole}. Its principal part gives the leading large-$d$ error. Theorem~\ref{zerofree} supplies a further real interval on which every section converges to $T$. These are local and sufficient statements: they do not require global continuation of one pole branch.

Near $q=1$, twisted Poisson summation has a different role. One Fourier frequency dominates for every root except $-1$, and its transform cancels in the quotient. This explains the independence of the odd limit from $k$. At $-1$, two equally large frequencies interfere and produce the real poles. We control the Fourier error and its derivative to prove simplicity and locate every sufficiently late pole. For $k=3$, Glasser's exact eta-cube transform \cite{Glasser,Coffey} yields an exponentially accurate elementary location formula. The transform is classical; its application to the zeros of the alternating denominator is the result used here.

The product identities are treated in \cite{Andrews,AndrewsAskeyRoy,GasperRahman}, the modular background in \cite{Apostol,Serre,DiamondShurman,KubertLang}, and Lambert-series rearrangements in \cite{Schmidt}. The eta transformation is also the cusp input in the partition asymptotics of Hardy--Ramanujan and Rademacher \cite{HardyRamanujan,Rademacher}. The CM targets are classical singular moduli \cite{Cox,Maier,Zagier}; here they serve as targets of explicit section approximation.

We use $X$ for the coefficient variable, $x=Xq^a$ for its rescaling, and $z$ for a root-of-unity sample. In analytic arguments $Q=q^a$ and, on the positive radius, $Q=e^{-t}$. The diagonal functions are $g_r$, their linear denominator is $L$, and the real Fourier kernel is $G$. The alternating denominator is $\mathcal D(Q)$, with $V_-(t)=Q\mathcal D(Q)$; $\varrho$ denotes the coefficient of the moving principal part. Formal $O(q^n)$ means divisibility by $q^n$. Analytic implied constants may depend on the fixed parameters; square roots and complex powers use principal branches.

\section{Formal contact and integral coefficients}
\begin{lemma}\label{lambert}
The numerator and denominator in \eqref{Rdef} satisfy
\begin{align}
 (1-X)U&=F(q)+X\sum_{n\ge1}f_n
 \frac{q^n(q^{an}-1)}{1-Xq^{an}},\label{LamU}\\
 (1-X)V&=F(q^a)+X\sum_{n\ge1}f_n
 \frac{q^{an}(q^{an}-1)}{1-Xq^{an}}.\label{LamV}
\end{align}
\end{lemma}
\begin{proof}
The constant coefficient in $X$ in the first identity is $F(q)$. For $m\ge1$, its coefficient on the left is
\[
 F(q^{am+1})-F(q^{a(m-1)+1})
 =\sum_{n\ge1}f_nq^{n(a(m-1)+1)}(q^{an}-1),
\]
which is the coefficient on the right after expanding the geometric denominator. The second identity follows in the same way, with $am+1$ replaced by $a(m+1)$. Each coefficient of a fixed power of $q$ is a finite sum.
\end{proof}

\begin{theorem}\label{contact}
Suppose $F(q)=1+f_pq^p+O(q^{p+1})$, where $p\ge1$ and $f_p\ne0$. For $m\ge1$,
\begin{equation}\label{Cfirst}
 C_m(q)=-f_p(1+f_p)^{m-1}q^{p(a(m-1)+1)}
       +O\!\left(q^{p(a(m-1)+1)+1}\right).
\end{equation}
For $d\ge2$,
\begin{equation}\label{Pfirst}
 P_d(q)-T(q)=-f_p(1+f_p)^{d-1}q^{(p+1)(a(d-1)+1)}
       +O\!\left(q^{(p+1)(a(d-1)+1)+1}\right).
\end{equation}
If $f_p\ne-1$, the displayed leading coefficients are nonzero.
\end{theorem}
\begin{proof}
Set $X=xq^{-ap}$ in the cleared forms \eqref{LamU}--\eqref{LamV}. Their lowest $q$-parts are
\[
 -\frac{f_px}{1-x}q^{-(a-1)p},\qquad
 \frac{1-(1+f_p)x}{1-x}.
\]
Indeed, a negative term in the $n$th numerator summand has the form $-f_nX^jq^{n(a(j-1)+1)}$, with $j\ge1$. After substitution its exponent is $aj(n-p)-(a-1)n$, at least $-(a-1)p$, with equality only for $n=p$. The positive terms have larger exponents. In the denominator the minimum is zero, again from $n=p$ and the constant term.

These computations take place in $K[[x]]((q))$. The second expression is a unit in $K[[x]]$, so the lowest part of their quotient is
\[
 -\frac{f_px}{1-(1+f_p)x}q^{-(a-1)p}
 =-f_p\sum_{m\ge1}(1+f_p)^{m-1}x^mq^{-(a-1)p}.
\]
This proves \eqref{Cfirst}. Since $C_0=F(q)/F(q^a)$,
\[
 P_d-T=q^{1-a}\sum_{j\ge1}C_{jd}q^{ajd}.
\]
The lower bound for the order of the $j$th summand is $(p+1)(a(jd-1)+1)$, which increases strictly with $j$. Thus $j=1$ gives \eqref{Pfirst}. Over a field the leading coefficient is nonzero when $f_p\ne0,-1$.
\end{proof}

\begin{corollary}\label{etacontact}
For $F=\Phi^k$ and $d\ge2$,
\begin{equation}\label{Wfirst}
 W_{a,d}(q,k)-T(q)
 =k(1-k)^{d-1}q^{2a(d-1)+2}
  +O\!\left(q^{2a(d-1)+3}\right).
\end{equation}
For a finite product $F=\prod_{\delta\ge1}\Phi(q^\delta)^{r_\delta}$ with integer exponents, let $p$ be the least index with $r_p\ne0$. The coefficient displayed in \eqref{Pfirst} is $r_p(1-r_p)^{d-1}$; it is the first nonzero coefficient when $r_p\ne1$.
\end{corollary}
\begin{proof}
The relevant first coefficients are $f_1=-k$ and $f_p=-r_p$, respectively. Substitute them into Theorem~\ref{contact}. The identities also hold polynomially in $k$.
\end{proof}

\begin{example}\label{mixedproduct}
For $F(q)=\Phi(q^2)^3\Phi(q^3)$, $a=3$, and $d=2$,
\[
 P_2(q)-T(q)=-6q^{12}+O(q^{13}).
\]
Indeed $F(q)=1-3q^2+O(q^3)$, so Theorem~\ref{contact} applies with $p=2$ and $f_p=-3$.
\end{example}

To examine the cancellation at $f_1=-1$, set $x=Xq^a$ and clear the first Lambert denominator. Define, within the formal sections,
\begin{align}
 N(x,q)={}&(1-x)F(q)+f_1x(q-q^{1-a})\notag\\
 &+x(1-x)\sum_{n\ge2}f_n
 \frac{q^{n-a}(q^{an}-1)}{1-xq^{a(n-1)}},\label{Ndef}\\
 M(x,q)={}&(1-x)F(q^a)+f_1x(q^a-1)\notag\\
 &+x(1-x)\sum_{n\ge2}f_n
 \frac{q^{a(n-1)}(q^{an}-1)}{1-xq^{a(n-1)}}.\label{Mdef}
\end{align}
Then $R(x/q^a,q)=N/M$. Both series have polynomial coefficients in $x$ at every $q$-degree, and $N$ has no powers below $q^{1-a}$.

\begin{theorem}\label{secondcontact}
Let $a\ge2$ and suppose, over a field, that
\[
 F(q)=1-q+f_s q^s+O(q^{s+1}),\qquad s\ge2,\qquad f_s\ne0.
\]
For $\ell\ge0$, the coefficients in \eqref{Rdef} satisfy
\begin{align}
 C_{2\ell+1}(q)&=(-f_s)^\ell q^{a(s+1)\ell+1}
                   +O(q^{a(s+1)\ell+2}),\label{Codd}\\
 C_{2\ell+2}(q)&=f_s(-f_s)^\ell q^{a(s+1)\ell+a+s}
                   +O(q^{a(s+1)\ell+a+s+1}).\label{Ceven}
\end{align}
Consequently, with $P_d$ defined by coefficient selection in \eqref{Pdef},
\begin{align}
 P_{2\ell+1}-T&=(-f_s)^\ell q^{a(s+3)\ell+2}
                   +O(q^{a(s+3)\ell+3})\quad(\ell\ge1),\label{Podd}\\
 P_{2\ell+2}-T&=f_s(-f_s)^\ell q^{a(s+3)\ell+2a+s+1}
                   +O(q^{a(s+3)\ell+2a+s+2})\quad(\ell\ge0).\label{Peven}
\end{align}
All displayed leading coefficients are nonzero. Over a commutative ring the coefficient identities remain valid, also when $f_s=0$, without the nonvanishing assertion.
\end{theorem}

The parity distinction comes from compositions of an integer: at an even
index, the least possible denominator order is attained only by parts of size
two. This selects different numerator terms at odd and even indices.
\begin{proof}
Put $Q=q^a$, $M_j=[x^j]M$ and $N_j=[x^j]N$, using
\eqref{Ndef}--\eqref{Mdef}. Direct expansion gives
\[
 M_0=1+O(Q),\qquad M_1=O(Q^{s-1}),\qquad
 M_j=f_sQ^{(s-1)(j-1)}+O(Q^{(s-1)(j-1)+1})\quad(j\ge2).
\]
Indeed, the $n$th summand contributing to $M_j$, for $j\ge2$, is
$f_n(Q^n-1)Q^{(n-1)(j-1)}(Q^{n-1}-1)$, whose first term occurs at $n=s$.
After division by $M_0$, expand $M^{-1}$ geometrically in its nonconstant
$x$-part. A composition $h=j_1+\cdots+j_v$ contributes at order at least
\[
 (s-1)\sum_{i=1}^v\max(1,j_i-1)\ge(s-1)\lceil h/2\rceil
\]
in $Q$. If $h=2\ell$, equality requires every part to be two. Hence
\begin{equation}\label{Minverse}
 \begin{aligned}
 [x^{2\ell}]M^{-1}
   &=(-f_s)^\ell Q^{(s-1)\ell}+O(Q^{(s-1)\ell+1}),\\
 [x^{2\ell+1}]M^{-1}&=O(Q^{(s-1)(\ell+1)}).
 \end{aligned}
\end{equation}
The numerator satisfies
\begin{align*}
 N_0&=F(q),\qquad N_1=q^{1-a}+O(q^{2-a}),\\
 N_j&=f_s q^{a(s-1)(j-2)+s-a}
       +O(q^{a(s-1)(j-2)+s-a+1})\quad(j\ge2).
\end{align*}
In $[x^{2\ell+1}](N/M)$, the unique least-order term is
$N_1[x^{2\ell}]M^{-1}$, of order $a(s-1)\ell+1-a$; all terms with
$j\ge2$ have order at least $a(s-1)\ell+s-a$, and the $j=0$ term has
still greater order. In $[x^{2\ell+2}](N/M)$ it is
$N_2[x^{2\ell}]M^{-1}$, of order $a(s-1)\ell+s-a$.
The $j=1$ term has greater order by at least $(a-1)(s-1)$, the $j=0$
term by at least $(a-1)s$, and every $j\ge3$ term by at least $a(s-1)$.
Multiplying by $q^{am}$ proves \eqref{Codd}--\eqref{Ceven}.
After adding the section weight $am+1-a$, these orders increase strictly
with $m$, so the first selected coefficient gives \eqref{Podd}--\eqref{Peven}.
Only series with constant coefficient one were inverted. Thus the coefficient
identities are valid over any commutative ring, and over a field the
hypothesis $f_s\ne0$ proves their nonvanishing.
\end{proof}

\begin{corollary}\label{PhiSecond}
For $F=\Phi$, the leading coefficient of $W_{a,d}(q,1)-T(q)$ is $(-1)^{d-1}$ and its order is
\[
 \begin{cases}
 5a(d-1)/2+2,&d\text{ odd},\\
 a(5d-6)/2+3,&d\text{ even}.
 \end{cases}
\]
In particular $W_{2,d}(q,1)-T(q)=(-1)^{d-1}q^{5d-3}+O(q^{5d-2})$.
\end{corollary}
\begin{proof}
Euler's product begins $\Phi(q)=1-q-q^2+O(q^3)$, so Theorem~\ref{secondcontact} applies with $s=2$ and $f_s=-1$.
\end{proof}

The cleared forms also control the coefficient diagonals and their integral divisibility.
Define the diagonal generating functions by
\begin{equation}\label{Gdef}
 R(x/q^a,q)=\sum_{r\ge0}q^{r-a+1}g_r(x),\qquad
 g_r(x)=\sum_{m\ge0}[q^{am-a+1+r}]C_m(q)x^m.
\end{equation}
The weight argument in Theorem~\ref{contact}, with weight $n-am$, proves that this substitution is well defined even when $f_1=0$. Put $L(x)=1-(1+f_1)x$.

\begin{lemma}\label{degreelemma}
Write $N=\sum_{w\ge1-a}q^wN_w(x)$ and $M=\sum_{i\ge0}q^{ai}M_{ai}(x)$ in \eqref{Ndef}--\eqref{Mdef}. Then
\[
 N_{1-a}=-f_1x,\qquad M_0=L,
\]
and
\begin{equation}\label{degreebounds}
 \deg M_{ai}\le i+1\quad(i\ge1),\qquad
 \deg N_{r-a+1}\le\lfloor(r-1)/a\rfloor+2\quad(r\ge0).
\end{equation}
\end{lemma}
\begin{proof}
Expand each denominator as a geometric series. In $N$, the negative part of the $n$th summand, $n\ge2$, has $q$-degree $n-a+a(n-1)j$ and $x$-degree at most $j+2$, where $j\ge0$. If its $q$-degree is $r-a+1$, then $n-1+a(n-1)j=r$. Hence $j\le\lfloor(r-1)/a\rfloor$. The positive part has an additional $an$ in its $q$-degree and satisfies the same bound. The terms outside the sum have degree at most one.

In $M$, the corresponding negative term has $q$-degree $a(n-1)(j+1)$ and degree at most $j+2$. If this is $ai$, then $j+2\le i+1$. The positive term again has higher $q$-degree. The constant $q$-parts follow directly from \eqref{Ndef}--\eqref{Mdef}.
\end{proof}

\begin{theorem}\label{diagonal}
For every $r\ge0$,
\begin{equation}\label{rationalG}
 g_r(x)=\frac{H_r(x)}{L(x)^{1+\lfloor r/a\rfloor}},\qquad H_r\in K[x],
\end{equation}
where
\begin{equation}\label{Hdegree}
 \deg H_r\le\lfloor r/a\rfloor+\lfloor(r-1)/a\rfloor+2.
\end{equation}
If $1+f_1\ne0$, the polynomial part of $g_r$ has degree at most $\lceil r/a\rceil$. For $m\ge\lceil r/a\rceil+1$, the coefficient $[x^m]g_r$ is $(1+f_1)^m$ times a polynomial in $m$ of degree at most $\lfloor r/a\rfloor$.
\end{theorem}
\begin{proof}
Comparing coefficients in $N=M\sum_{r\ge0}q^{r-a+1}g_r$ gives
\[
 Lg_r=N_{r-a+1}-\sum_{i=1}^{\lfloor r/a\rfloor}M_{ai}g_{r-ai}.
\]
Multiplying by $L^{\lfloor r/a\rfloor}$ yields the finite recursion
\begin{equation}\label{Hrec}
 H_r=L^{\lfloor r/a\rfloor}N_{r-a+1}
 -\sum_{i=1}^{\lfloor r/a\rfloor}L^{i-1}M_{ai}H_{r-ai}.
\end{equation}
The initial numerator is $H_0=-f_1x$. Lemma~\ref{degreelemma} and induction prove that every $H_r$ is a polynomial. The first term in \eqref{Hrec} satisfies \eqref{Hdegree}; for the $i$th term in the sum, the degree is at most
\[
 (i-1)+(i+1)+\lfloor r/a\rfloor-i+\lfloor(r-1)/a\rfloor-i+2,
\]
which gives the same bound.

When $L$ has degree one, polynomial division shows that the polynomial part has degree at most $\lfloor(r-1)/a\rfloor+1=\lceil r/a\rceil$. The remaining proper fraction is a linear combination of $L^{-j}$, $1\le j\le1+\lfloor r/a\rfloor$. Since
\[
 [x^m]L^{-j}=\binom{m+j-1}{j-1}(1+f_1)^m,
\]
the final assertion follows. This is the usual coefficient description for rational generating functions; compare \cite[Chapter 4]{Stanley} and \cite[Chapter IV]{FlajoletSedgewick}.
\end{proof}

The increasing denominator powers in \eqref{rationalG} can be collected into one formal pole. Theorem~\ref{analyticpole} will identify this formal root with the unique nearest analytic pole when $F=\Phi^k$, $k\ge2$, and $|q|$ is sufficiently small.

\begin{theorem}\label{movingpole}
Suppose $f_1\ne0,-1$. There is a unique $x_0(q)\in K[[q^a]]$ with
\[
 x_0(0)=\frac1{1+f_1},\qquad M(x_0(q),q)=0.
\]
Set
\begin{equation}\label{rho}
 \varrho(q)=-\frac{N(x_0(q),q)}{M_x(x_0(q),q)}
 =-\frac{f_1}{(1+f_1)^2}q^{1-a}+O(q^{2-a}).
\end{equation}
For every $r\ge0$, the coefficient of $q^{r-a+1}$ in
\begin{equation}\label{principalpart}
 R(x/q^a,q)-\frac{\varrho(q)}{x_0(q)-x}
\end{equation}
is a polynomial in $x$ of degree at most $\lceil r/a\rceil$. Therefore
\begin{equation}\label{stable}
 [q^{am-a+1+r}]C_m(q)
 =[q^{r-a+1}]\varrho(q)x_0(q)^{-m-1}
 \quad\text{if }m\ge\lceil r/a\rceil+1.
\end{equation}
\end{theorem}
\begin{proof}
The constant term of $M$ in $q$ is $L$, whose derivative is the nonzero scalar $-(1+f_1)$. Equating successive coefficients in $M(x_0(q),q)=0$ determines a unique root: at each order the unknown coefficient is multiplied by this scalar. Since $M$ involves only $q^a$, the root lies in $K[[q^a]]$.

All substitutions are defined because the coefficients of $M$ and $N$ at fixed $q$-degree are polynomials. The difference quotient
\[
 \widehat M(x,q)=\frac{M(x,q)-M(x_0(q),q)}{x-x_0(q)}
\]
belongs to $K[x][[q]]$ and has nonzero scalar constant term $-(1+f_1)$. Thus $N/\widehat M$ belongs to $K[x]((q))$. Denote this quotient by $\mathcal B(x,q)$. Polynomial division at each $q$-degree gives
\[
 \frac{N}{M}
 =\frac{\mathcal B(x,q)-\mathcal B(x_0(q),q)}{x-x_0(q)}
   +\frac{\varrho(q)}{x_0(q)-x},
 \qquad \varrho(q)=-\mathcal B(x_0(q),q).
\]
The first term has polynomial coefficients in $x$. If $x_0(q)=x_0(0)+\delta(q)$, the second term expands as
\[
 \varrho(q)\sum_{j\ge0}\frac{(-\delta(q))^j}{(x_0(0)-x)^{j+1}}.
\]
At each $q$-degree this is a proper rational function of $x$. It follows that the first term gives exactly the polynomial part of $g_r$, whose degree was bounded in Theorem~\ref{diagonal}. Expanding $(x_0-x)^{-1}=\sum_{m\ge0}x^mx_0^{-m-1}$ proves \eqref{stable}. Finally $\widehat M(x_0,q)=M_x(x_0,q)$ and $N_{1-a}=-f_1x$ give \eqref{rho}.
\end{proof}

\begin{corollary}\label{kpolynomial}
For $a=2$ and $F=\Phi^k$, let $\gamma_{m,r}=[q^{2m-1+r}]C_m(q)$. For each $r\ge0$ there is a polynomial $P_r(m;k)\in\Q[k,m]$, of degree at most $\lfloor r/2\rfloor$ in $m$, such that
\begin{equation}\label{kfactor}
 \gamma_{m,r}=(-1)^{m-1}k(k-1)^{m-1-r}P_r(m;k)
 \qquad(m\ge r+1).
\end{equation}
For each integer $m\ge r+1$, $P_r(m;1)\ne0$. In particular the exact power of $k-1$ dividing $\gamma_{m,r}\in\Q[k]$ is $m-1-r$.
\end{corollary}
\begin{proof}
Every nonconstant coefficient of $\Phi^k$ is divisible by $k$ in $\Q[k]$. Reduction of \eqref{Rdef} modulo $k$ gives $R=1$. With $s=1+\lfloor r/2\rfloor$ and $L=1+(k-1)x$, Theorem~\ref{diagonal} therefore gives
\[
 g_r(x)-g_r(0)=\frac{kx\Pi_r(x,k)}{L^s},\qquad
 \Pi_r=\sum_{i=0}^r\pi_i(k)x^i\in\Q[k,x].
\]
For $m\ge r+1$, coefficient extraction proves \eqref{kfactor} with
\[
 P_r(m;k)=\sum_{i=0}^r(-1)^i\pi_i(k)(k-1)^{r-i}
                 \binom{m-i+s-2}{s-1}.
\]
Its degree in $m$ is at most $s-1$. At $k=1$, Corollary~\ref{PhiSecond} gives
$[x^{r+1}]g_r(x)=(-1)^r$. Since $L=1$ at this specialization, $\pi_r(1)=(-1)^r$. Thus
\[
 P_r(m;1)=\binom{m-r+s-2}{s-1}>0.
\]
The same calculation includes $r=0$, where $P_0=1$.
\end{proof}

\begin{example}
For $a=2$, $F=\Phi^k$, and $L=1+(k-1)x$, direct extraction from \eqref{Hrec} gives
\[
 g_0=\frac{kx}{L},\qquad
 g_1=1+\frac{kx(kx-k-3x+1)}{2L}.
\]
Hence $\gamma_{m,0}=(-1)^{m-1}k(k-1)^{m-1}$ and, for $m\ge2$,
\[
 \gamma_{m,1}=(-1)^m\frac{k(k-2)(k+1)}2(k-1)^{m-2}.
\]
These identities also follow by using $f_1=-k$ and $f_2=k(k-3)/2$ in \eqref{Ndef}--\eqref{Mdef}.
\end{example}

The degree bound also gives an integral statement. It is important here to use the integral recursion \eqref{Hrec}, rather than infer divisibility of integer values from a factorization over $\Q[k]$.

\begin{theorem}\label{adic}
Let $a,d\ge2$ and $k\in\Z$ with $h=|k-1|\ge2$. Define
\[
 E_0=2a(d-1)+2,\qquad
 E_d=\begin{cases}
 5a(d-1)/2+2,&d\text{ odd},\\
 a(5d-6)/2+3,&d\text{ even}.
 \end{cases}
\]
The coefficients of $W_{a,d}(q,k)-T(q)$ below $q^{E_0}$ vanish. For $0\le r\le E_d-E_0$,
\begin{equation}\label{generaladic}
 h^{\,d-2-\lfloor r/a\rfloor-\lfloor(r-1)/a\rfloor}
 \ \bigm|\ [q^{E_0+r}]\bigl(W_{a,d}(q,k)-T(q)\bigr).
\end{equation}
The exponent in \eqref{generaladic} is nonnegative. At the endpoint, with $b=-k(k-3)/2$,
\begin{equation}\label{adicendpoint}
 [q^{E_d}]\bigl(W_{a,d}(q,k)-T(q)\bigr)
 \equiv(-1)^{d-1}b^{\lfloor d/2\rfloor}\pmod h.
\end{equation}
If $h\ge3$, this endpoint coefficient is nonzero modulo $h$, and therefore
is the first nonzero coefficient modulo $h$. More explicitly,
\[
 b\equiv
 \begin{cases}
 1\pmod h,&h\text{ odd},\\
 1+h/2\pmod h,&h\text{ even}.
 \end{cases}
\]
\end{theorem}
\begin{proof}
All coefficients of $\Phi^k$ are integers, including for negative integer $k$. The recursion \eqref{Hrec} therefore gives $H_r\in\Z[x]$ after specialization, with the degree bound \eqref{Hdegree}. Put $s=1+\lfloor r/a\rfloor$ and write $H_r=\sum_j u_jx^j$. Since $L=1+(k-1)x$,
\[
 [x^d]g_r(x)=\sum_{0\le j\le d}
 u_j(-1)^{d-j}(k-1)^{d-j}
       \binom{d-j+s-1}{s-1}.
\]
For the stated range of $r$, the degree bound is at most $d$, so every summand is divisible by
$h^{d-2-\lfloor r/a\rfloor-\lfloor(r-1)/a\rfloor}$.

By Theorem~\ref{contact}, the term selected from $C_{2d}$ starts at $q^{4ad-2a+2}$, and later selected terms start still later. Since $E_d<4ad-2a+2$, only $C_d$ contributes through $q^{E_d}$. Its coefficient at $q^{E_0+r}$ after selection is precisely $[x^d]g_r$. This proves \eqref{generaladic} and the vanishing below $E_0$. Directly from the two definitions of $E_d$, the exponent in \eqref{generaladic} is positive for $r<E_d-E_0$ and zero at the endpoint.

Reduce the cleared forms modulo $h$. Then $f_1=-k=-1$ and $f_2=k(k-3)/2=-b$ in $\Z/h\Z$. The coefficient computation in the proof of Theorem~\ref{secondcontact} uses no division by $f_2$ and is valid over this ring even if $f_2$ is zero or a zero divisor. For $d=2\ell+1$ it gives the endpoint $b^\ell$, and for $d=2\ell+2$ it gives $-b^{\ell+1}$. This is \eqref{adicendpoint}. For odd $h$, reduction of $2b=-k(k-3)$ modulo $h$ gives $b=1$.
For even $h$, write $k=1+\epsilon h$, with $\epsilon\in\{1,-1\}$. Then
\[
 b=1+\frac{\epsilon h}{2}-\frac{h^2}{2}
   \equiv1+\frac h2\pmod h.
\]
If $4\mid h$, this residue has square one. If $h\equiv2\pmod4$, it is
idempotent. In either case, for $h>2$ every positive power is nonzero
modulo $h$. This proves the endpoint assertion.
\end{proof}

\begin{corollary}\label{leveltwoadic}
For $a=2$, integer $k$ with $|k-1|\ge2$, and $4d-2\le n\le5d-3$,
\[
 (k-1)^{5d-3-n}\ \bigm|\ [q^n]\bigl(W_{2,d}(q,k)-T(q)\bigr).
\]
For $k=24$ the defect modulo $23$ begins exactly with
\[
 W_{2,d}(q,24)-T(q)\equiv(-1)^{d-1}q^{5d-3}+O(q^{5d-2})\pmod{23}.
\]
\end{corollary}
\begin{proof}
For $a=2$, $\lfloor r/2\rfloor+\lfloor(r-1)/2\rfloor=r-1$ and $E_d=5d-3$. Apply Theorem~\ref{adic}; its last assertion applies to $h=23$.
\end{proof}

\begin{corollary}\label{modtwoendpoint}
For $k=3$ or $k=-1$ and $a,d\ge2$, the first nonzero coefficient of
$W_{a,d}(q,k)-T(q)$ modulo two is one and occurs at
\[
 \begin{cases}
 3a(d-1)+2,&d\text{ odd},\\
 a(3d-4)+4,&d\text{ even}.
 \end{cases}
\]
In particular,
\[
 W_{2,d}(q,k)-T(q)\equiv q^{6d-4}+O(q^{6d-3})\pmod2.
\]
\end{corollary}
\begin{proof}
The coefficients of the Euler products give
\[
 \Phi(q)^3=1-3q+5q^3+O(q^4),\qquad
 \Phi(q)^{-1}=1+q+2q^2+3q^3+O(q^4).
\]
Both reduce to $1+q+q^3+O(q^4)$ over $\mathbb F_2$. Theorem~\ref{secondcontact}
therefore applies with $s=3$ and $f_s=1$. Substituting
$d=2\ell+1$ and $d=2\ell+2$ into \eqref{Podd}--\eqref{Peven} proves the result.
\end{proof}

\begin{example}\label{evenmodulus}
For $k=5$, one has $h=4$ and $b=-5\equiv3\pmod4$; the endpoint is nonzero
for every $d$, and for $d=2$ it equals one modulo four. At the exceptional
modulus two,
\[
 W_{2,2}(q,3)-T(q)=-6q^6+6q^7-33q^8+O(q^9).
\]
Thus the first nonzero term modulo two is $q^8$, in agreement with
Corollary~\ref{modtwoendpoint}.
\end{example}

\section{Analytic poles and section convergence}
Fix $k\in\Z_{>0}$. For $0<|q|<1$, define
\begin{equation}\label{Jdef}
 J(z,q)=\frac{\sum_{m\ge0}z^mq^{am+1}\Phi(q^{am+1})^k}
              {\sum_{m\ge0}z^mq^{a(m+1)}\Phi(q^{a(m+1)})^k}.
\end{equation}
Both sums converge locally uniformly for $|z|<|q|^{-a}$. Their quotient is meromorphic, and the denominator is not identically zero. Near $q=0$, $J(z,q)=q^{1-a}R(zq^a,q)$. Hence
\begin{equation}\label{Wanalytic}
 W_{a,d}(q,k)=\frac1d\sum_{z^d=1}J(z,q)
\end{equation}
gives a meromorphic continuation of the formal section to $0<|q|<1$. 
The sampling and coefficient variables are related by
\[
 Q=q^a,\qquad X=zQ,\qquad x=XQ=zQ^2,\qquad z_0=x_0/Q^2.
\]
The first theorem concerns poles in $z$ with $q$ fixed. The radial theorem later concerns poles in $q$ at the fixed sample $z=-1$. The denominator $M$ in \eqref{Mdef} depends on $q$ only through $Q$.

\begin{theorem}\label{analyticpole}
Fix integers $a\ge2$, $k\ge2$, and a real number $c>1/(k-1)$. For all sufficiently small nonzero $q$, the formal root in Theorem~\ref{movingpole} converges to a function analytic in $Q$ near zero. The meromorphic continuation of $z\mapsto J(z,q)$ has exactly one pole in $|z|\le c/|Q|^2$. It is simple, with
\begin{equation}\label{analyticlocation}
 z_0(q)=\frac{x_0(q)}{Q^2}
       =-\frac{1+O(Q)}{(k-1)Q^2},\qquad
 \Res_{z=z_0}J(z,q)=-\frac{q^{1-a}\varrho(q)}{Q^2}.
\end{equation}
For fixed positive $q$ sufficiently small, $z_0<-1$ and, for any fixed $r_0$ with $|z_0|<r_0<c/Q^2$,
\begin{equation}\label{geometricerror}
 W_{a,d}(q,k)-T(q)
 =\frac{q^{1-a}\varrho(q)}{x_0(q)}
          \frac{z_0(q)^{-d}}{1-z_0(q)^{-d}}+O_{q,r_0}(r_0^{-d})
 \quad(d\to\infty).
\end{equation}
The leading error has sign $(-1)^{d-1}$ for all sufficiently large $d$.
\end{theorem}
\begin{proof}
On every fixed disk $|x|\le c$, the series \eqref{Mdef} converges normally for small $|Q|$. Indeed its denominators with $n\ge2$ are bounded away from zero once $c|Q|<1/2$, and convergence follows from the convergence of $F$ in the unit disk. The same argument shows that $q^{a-1}N$ is jointly holomorphic for small $q$ and bounded $x$. Uniformly on this disk,
\[
 M(x,q)=1+(k-1)x+O(Q),\qquad
 q^{a-1}N(x,q)=kx+O(q).
\]
On $|x|=c$, the first limiting polynomial has modulus at least $(k-1)c-1>0$. Rouch\'e's theorem therefore gives exactly one zero of $M$ in the disk, counted with multiplicity. The implicit function theorem at $(x,Q)=(-1/(k-1),0)$ makes this zero simple and analytic in $Q$. Uniqueness of the formal solution identifies its Taylor series with Theorem~\ref{movingpole}.

At this zero, $q^{a-1}N\to-k/(k-1)\ne0$, so there is no numerator cancellation. The cleared quotient is a meromorphic continuation of $R(x/Q,q)$. Clearing the first Lambert factor introduces no singularity at $x=1$, since $M(1,q)=k(1-Q)\ne0$. Thus $J(z,q)=q^{1-a}N(zQ^2,q)/M(zQ^2,q)$ has precisely the asserted pole. Its residue is the negative of the coefficient of $(x_0-x)^{-1}$, divided by the scale factor $Q^2$, which proves \eqref{analyticlocation}.

For positive small $q$, reality and uniqueness give a real root with $x_0<0$, while \eqref{rho} gives $\varrho(q)>0$. Also $|z_0|>1$. The function
\[
 H(z)=J(z,q)-\frac{q^{1-a}\varrho(q)}{x_0(q)}
                         \frac1{1-z/z_0(q)}
\]
is holomorphic on $|z|\le r_0$. Cauchy's estimate and the root filter give
$d^{-1}\sum_{z^d=1}H(z)-H(0)=O_{q,r_0}(r_0^{-d})$. Averaging the geometric fraction gives \eqref{geometricerror}. Its prefactor is negative, $z_0$ is negative, and $r_0>|z_0|$, so the remainder is smaller than the displayed nonzero leading term. The sign assertion follows.
\end{proof}

\begin{remark}\label{poleconsistency}
For $k\ge2$, the pole term in \eqref{geometricerror} also recovers the formal first defect, with $d$ fixed. By Theorem~\ref{movingpole},
\[
 \varrho=\frac{kq^{1-a}}{(k-1)^2}(1+O(q)),\qquad
 x_0=-\frac{1+O(Q)}{k-1}.
\]
Consequently its expansion at $q=0$ is
\begin{equation}\label{formalprincipal}
 \frac{q^{1-a}\varrho}{x_0}
 \frac{(Q^2/x_0)^d}{1-(Q^2/x_0)^d}
 =k(1-k)^{d-1}q^{2a(d-1)+2}(1+O_d(q)).
\end{equation}
This comparison does not require uniformity of the analytic remainder as $q\to0$. Indeed the coefficient of $q^{r-a+1}$ in
$R(x/Q,q)-\varrho/(x_0-x)$ is a polynomial of degree at most $\lceil r/a\rceil$, by Theorem~\ref{movingpole}. A nonconstant term surviving the root filter has degree $m\ge d$, hence $r\ge a(m-1)+1$. Substituting $x=zQ^2$ and multiplying by $q^{1-a}$ gives exponent at least $3a(m-1)+3$. Thus the difference between $W_{a,d}-T$ and the left side of \eqref{formalprincipal} is formally $O(q^{3a(d-1)+3})$, which proves consistency with \eqref{Wfirst} independently of the large-$d$ limit.
\end{remark}


There is also a real interval that can be obtained without following a zero branch. The coefficient argument is of Enestr\"om--Kakeya type; see \cite{Melman} for the classical polynomial setting and related zero-exclusion methods. For $0<Q<1$, set
\begin{equation}\label{threshold}
 \beta_k(Q)=Q\prod_{n\ge1}(1+Q^n)^k,
 \qquad \beta_k(Q_*(k))=1.
\end{equation}
The function $\beta_k$ is strictly increasing from zero to infinity, so $Q_*(k)$ exists and is unique. The limit at $1$ follows, for example, by retaining arbitrarily many factors in the product.

\begin{theorem}\label{zerofree}
For integers $a\ge2$, $k\ge1$ and real $0<q<1$, put $Q=q^a$. The denominator of $J(z,q)$ has no zero on
\[
 |z|\le\beta_k(Q)^{-1}.
\]
Consequently every section has no pole for $0<q\le Q_*(k)^{1/a}$, and $W_{a,d}(q,k)\to T(q)$ exponentially as $d\to\infty$ at each such fixed nome. If $Q<Q_*(k)$, the error is $O_q(\beta_k(Q)^d)$.
\end{theorem}
\begin{proof}
Apart from the nonzero factor $Q$, the denominator is $\sum_{m\ge0}b_mz^m$, where $b_m=Q^m\Phi(Q^{m+1})^k>0$. Its successive ratios satisfy
\[
 \frac{b_{m+1}}{b_m}
 =Q\prod_{n\ge1}
 \left(\frac{1-Q^{(m+2)n}}{1-Q^{(m+1)n}}\right)^k.
\]
For each $n$, the function $(1-Q^nu)/(1-u)$ is strictly increasing for $0<u<1$. Taking $u=Q^{(m+1)n}$ shows that the ratios are strictly decreasing in $m$. The first ratio is $\beta_k(Q)$, and their limit is $Q$.

Put $r_0=\beta_k(Q)^{-1}$ and $c_m=b_mr_0^m$. Then $c_0=c_1>c_2>\cdots>0$ and $c_m\to0$ geometrically, since $Qr_0<1$. For $S(w)=\sum_{m\ge0}c_mw^m$,
\[
 (1-w)S(w)=c_0-\sum_{m\ge1}(c_{m-1}-c_m)w^m,
 \qquad \sum_{m\ge1}(c_{m-1}-c_m)=c_0.
\]
For $|w|<1$, the second sum has modulus strictly less than $c_0$. On $|w|=1$, equality sufficient for cancellation would require $w^m=1$ for every $m\ge2$, since all those differences are positive. The conditions for $m=2,3$ force $w=1$, where $S(1)>0$. Thus $S$ is nonzero on the closed unit disk, proving the stated zero-free disk in $z$.

If $Q\le Q_*(k)$, this disk contains the unit disk. Both sums defining $J$ are holomorphic in $|z|<Q^{-1}$, and $r_0<Q^{-1}$. Compactness therefore gives a slightly larger zero-free disk. Cauchy's estimate and the root filter prove exponential convergence. When $Q<Q_*(k)$, use the circle $|z|=r_0>1$ directly: the error is bounded by a constant times $r_0^{-d}/(1-r_0^{-d})$, proving the final estimate.
\end{proof}

\begin{corollary}\label{modularthreshold}
For $0<Q<1$, write $Q=e^{2\pi i\tau}$ with $\tau$ positive imaginary and put $u(\tau)=\Delta(\tau)/\Delta(2\tau)$. Then
\begin{equation}\label{betamodular}
 \beta_{24}(Q)=u(\tau)^{-1},\qquad
 \beta_k(Q)=Q^{1-k/24}u(\tau)^{-k/24}\quad(k\ge1),
\end{equation}
with positive real powers. The threshold $Q_*(24)$ is the unique positive nome for which $\Delta(\tau_*)=\Delta(2\tau_*)$, and
\begin{equation}\label{thresholdj}
 j(\tau_*)=257^3,\qquad j(2\tau_*)=17^3,
\end{equation}
where $j$ is the classical modular invariant. For $k=24$, the zero-free interval in Theorem~\ref{zerofree} is equivalently $u(\tau)\ge1$, with $Q=q^a$.
\end{corollary}
\begin{proof}
The product definition gives
\[
 u(\tau)=Q^{-1}\left(\frac{\Phi(Q)}{\Phi(Q^2)}\right)^{24}
        =Q^{-1}\prod_{n\ge1}(1+Q^n)^{-24},
\]
proving \eqref{betamodular}. This function decreases strictly from infinity to zero as $Q$ increases from zero to one. Hence $u=1$ characterizes the threshold and $u\ge1$ characterizes the stated interval. The classical relation
$j(\tau)=(u+256)^3/u^2$, recalled in \eqref{jrelation}, gives the first value in \eqref{thresholdj}. The eta transformation gives
$u(-1/(2\tau))=4096/u(\tau)$. Applying the same relation there and using $j(-1/(2\tau))=j(2\tau)$ yields
$j(2\tau)=(u+16)^3/u$, proving the second value.
\end{proof}

The condition $j(\tau)=257^3$ alone does not specify this point on the whole positive imaginary axis: it also holds at $-1/\tau_*$. The condition $u=1$ selects the required representative. These are algebraic modular values; no algebraicity or CM assertion about $\tau_*$ is used here.

For later use, define the single alternating denominator
\begin{equation}\label{commonD}
 \mathcal D(Q)=\sum_{m\ge0}(-1)^m Q^m\Phi(Q^{m+1})^k.
\end{equation}
For every dilation, the denominator of $J(-1,q)$ is $Q\mathcal D(Q)$. Thus its possible real poles share the same denominator zeros in the variable $Q$. Theorem~\ref{poles} proves that all sufficiently late such zeros are simple and are actual poles of every fixed even section. Theorem~\ref{zerofree} excludes all real section poles for $Q\le Q_*(k)$. Neither theorem asserts that every denominator zero is uncancelled or that one local zero branch remains simple throughout $0<Q<1$.

For a general fixed nome, we describe the fixed-nome limit when the quotient has interior poles. This is the contour interpretation of the exponentially convergent trapezoidal rule \cite{TrefethenWeideman}.

\begin{proposition}\label{contourlimit}
Let $a\ge2$ and $k\ge1$ be integers, take $F=\Phi^k$, and fix $0<|q|<1$. Suppose $J(z,q)$ has no pole on $|z|=1$. Then
\begin{equation}\label{contourformula}
 \lim_{d\to\infty}W_{a,d}(q,k)
 =\frac1{2\pi i}\int_{|z|=1}\frac{J(z,q)}z\dd z
 =T(q)+\sum_{|z_\nu|<1}\Res_{z=z_\nu}\frac{J(z,q)}z,
\end{equation}
where the sum is over the actual poles of $J$. The convergence is exponential in $d$.
\end{proposition}
\begin{proof}
The quotient is meromorphic in $|z|<|q|^{-a}$, a disk containing the closed unit disk. Its poles are discrete, and by hypothesis none lies on the unit circle. Hence it is holomorphic in an annulus $r<|z|<s$ with $r<1<s$. Write its Laurent expansion there as $J(z,q)=\sum_{n\in\Z}b_nz^n$. On slightly smaller bounding circles, Cauchy's estimates give
\[
 |b_n|\le C s_0^{-n}\quad(n\ge0),\qquad
 |b_{-n}|\le C r_0^n\quad(n\ge1),
 \qquad r<r_0<1<s_0<s.
\]
The expansion converges absolutely on $|z|=1$, so averaging over the $d$th roots gives
\[
 W_{a,d}(q,k)=\sum_{j\in\Z}b_{jd}
 =b_0+O(r_0^d+s_0^{-d}).
\]
The contour integral equals $b_0$. The residue theorem gives the last equality in \eqref{contourformula}, since $J(0,q)=T(q)$ and there are only finitely many poles inside the unit circle.
\end{proof}

Theorems~\ref{analyticpole}, \ref{zerofree}, and \ref{CMtheorem} give proved regimes in which the residue sum is empty. The first also identifies the leading geometric error at a fixed sufficiently small nome. Proposition~\ref{contourlimit} explains the additional residues at a general nome. These analytic distinctions coexist with the exact formal contact of Theorem~\ref{contact}.


\section{Radial limits and pole geometry}
Fix $k\in\Z_{>0}$. In this section write $q=e^{-t/a}$, with $t>0$, and put
\begin{equation}\label{Ganalytic}
 G(v)=e^{-v}\Phi(e^{-v})^k\quad(v>0),\qquad G(v)=0\quad(v\le0),
 \qquad \widehat G(\xi)=\int_{\mathbb R}G(v)e^{-2\pi i\xi v}\dd v.
\end{equation}

\begin{lemma}\label{Gschwartz}
For $\Re v>0$,
\begin{equation}\label{Gtransform}
 G(v)=\left(\frac{2\pi}{v}\right)^{k/2}
 \exp\!\left[-\left(1-\frac{k}{24}\right)v-\frac{\pi^2k}{6v}\right]
 \Phi(e^{-4\pi^2/v})^k.
\end{equation}
The real function in \eqref{Ganalytic} is a Schwartz function, and $\widehat G(0)>0$.
\end{lemma}
\begin{proof}
With $\tau=iv/(2\pi)$, the classical transformation
$\eta(-1/\tau)=(-i\tau)^{1/2}\eta(\tau)$ gives
\[
 \Phi(e^{-v})=\left(\frac{2\pi}{v}\right)^{1/2}
 e^{v/24-\pi^2/(6v)}\Phi(e^{-4\pi^2/v}).
\]
The identity holds for positive $v$ and then throughout $\Re v>0$ by analytic continuation, with principal powers. Raising to the $k$th power proves \eqref{Gtransform}. The transformation and its cusp consequences are classical; see \cite{Apostol,Rademacher}.

As $v\downarrow0$, the factor $e^{-\pi^2k/(6v)}$ dominates every power of $v^{-1}$. The remaining product, and each of its derivatives after multiplication by a suitable power of $v$, are bounded. Thus $G$ and all its right derivatives vanish at zero. As $v\to+\infty$, the defining product shows exponential decay of every derivative. The zero extension is therefore Schwartz. Positivity on $(0,\infty)$ gives $\widehat G(0)>0$. This is the Laplace transform of $\Phi(e^{-v})^k$ at $1$; for $k=3$, Lemma~\ref{cubictransform} gives $2\pi/\cosh(\pi\sqrt7/2)$.
\end{proof}

The derivative estimate in the next lemma is needed to distinguish simple poles from mere sign changes.

\begin{lemma}\label{Fourier}
Put
\begin{equation}\label{Bk}
 B_k=\sqrt\pi\,(2\pi)^{(3k-3)/4}
                \left(\frac{\pi^2k}{6}\right)^{(1-k)/4}>0.
\end{equation}
As $\xi\to+\infty$,
\begin{equation}\label{Fourierexp}
 \widehat G(\xi)=B_k\xi^{(k-3)/4}
 \exp\!\left[-(1+i)\pi^{3/2}\sqrt{\frac{2k\xi}{3}}
                 +\frac{(k-3)\pi i}{8}\right](1+\varepsilon(\xi)),
\end{equation}
where
\begin{equation}\label{Fourierderror}
 \varepsilon(\xi)=O(\xi^{-1/2}),\qquad
 \varepsilon'(\xi)=O(\xi^{-3/2}).
\end{equation}
The first expansion holds uniformly in some fixed complex sector about the positive real axis, for the analytic continuation defined by a rotated integral.
\end{lemma}
\begin{proof}
For this proof put $A=\pi^2k/6$, $c=1-k/24$, and $b=c+2\pi i\xi$. Rotate the path in \eqref{Ganalytic} to $v=re^{-i\pi/4}$, $r>0$. At the origin, \eqref{Gtransform} bounds the integrand by a power of $r^{-1}$ times $e^{-A/(\sqrt2r)}$. At infinity, the original product gives exponential decay throughout the closed sector $-\pi/4\le\arg v\le0$, while $|e^{-2\pi i\xi v}|\le1$ for real $\xi>0$. The connecting arcs therefore vanish.

On the rotated path, replace the final product in \eqref{Gtransform} by $1$. The resulting integral is
\begin{equation}\label{I0}
 I_0(\xi)=2(2\pi)^{k/2}
       (A/b)^{(2-k)/4}K_{k/2-1}(2\sqrt{Ab}).
\end{equation}
For positive real $b$ this is the usual integral representation of $K$, and continuation gives it in the sector $-\pi/4<\arg b<3\pi/4$, where $\Re(be^{-i\pi/4})>0$ and the rotated integral converges. This condition holds for all sufficiently large positive $\xi$, also when $k>24$ and $c<0$. The large-argument expansion
\[
 K_\nu(z)=\sqrt{\frac\pi{2z}}e^{-z}(1+O(z^{-1}))
\]
is uniform on closed subsectors containing the arguments used here; see \cite[\S\S10.32, 10.40]{DLMF}. Substitution into \eqref{I0} gives
\[
 I_0(\xi)=(2\pi)^{k/2}\sqrt\pi\,
 A^{(1-k)/4}b^{(k-3)/4}e^{-2\sqrt{Ab}}(1+O(|\xi|^{-1/2})).
\]
Since $b=2\pi i\xi(1+O(\xi^{-1}))$, this is \eqref{Fourierexp} with the constant \eqref{Bk}, up to the product remainder. In particular $B_3=4\pi$.

Let $\omega(v)=\Phi(e^{-4\pi^2/v})^k-1$. On the rotated path and for $0<r\le1$, the elementary product bound gives
\[
 |\omega(re^{-i\pi/4})|\le C_k e^{-4\pi^2/(\sqrt2r)}.
\]
For $r\ge1$, it gives $|\omega(re^{-i\pi/4})|\le C_ke^{C_kr}$. To see both estimates, put $u=e^{-4\pi^2/(\sqrt2r)}$ and use
\[
 |\omega(v)|\le\prod_{n\ge1}(1+u^n)^k-1
            \le\exp\!\left(\frac{ku}{1-u}\right)-1.
\]
For large complex $\xi$ in a sufficiently small sector about the positive axis, $\beta=\Re(be^{-i\pi/4})$ is positive and comparable to $|\xi|$. The part of $\widehat G-I_0$ with $r\le1$ is bounded by
\begin{equation}\label{Fouriertail}
 C_k\int_0^\infty r^{-k/2}
       e^{-(A+4\pi^2)/(\sqrt2r)-\beta r}\dd r.
\end{equation}
The positive-real $K$ integral bounds \eqref{Fouriertail} by a power of $|\xi|$ times
\[
 \exp\!\left[-2\sqrt{(A+4\pi^2)\beta/\sqrt2}\right].
\]
On the positive real axis its exponential rate is
$2\sqrt{\pi(A+4\pi^2)\xi}$, strictly larger than the rate $2\sqrt{\pi A\xi}$ of $I_0$. The two rates depend continuously on $\arg\xi$. Shrinking the sector preserves a fixed positive gap, so \eqref{Fouriertail} divided by the modulus of the leading term in \eqref{Fourierexp} is $O(e^{-c_k\sqrt{|\xi|}})$. For $r\ge1$, the above bound for $\omega$ gives $O(e^{-c_k|\xi|})$.

The rotated integral defines an analytic continuation in this sector. We have proved that its relative error in \eqref{Fourierexp} is analytic and uniformly $O(|\xi|^{-1/2})$. For large real $\xi$, a disk of radius $\delta\xi$ lies in a slightly larger admissible sector. Cauchy's derivative estimate on this disk gives $\varepsilon'(\xi)=O(\xi^{-3/2})$. This proves \eqref{Fourierderror}, with the uniform remainder control used in classical asymptotic analysis \cite{Olver,Wong}.
\end{proof}

\begin{lemma}\label{twists}
For real $\alpha$ and $0<c\le1$, let
\[
 S(c,\alpha;t)=\sum_{m\in\Z}e^{2\pi i\alpha m}G(t(m+c)).
\]
Then
\begin{equation}\label{Poisson}
 S(c,\alpha;t)=\frac1t\sum_{n\in\Z}
 e^{2\pi i(n-\alpha)c}\widehat G((n-\alpha)/t).
\end{equation}
For each fixed $-1/2<\alpha<1/2$, the denominator of $J(e^{2\pi i\alpha},e^{-t/a})$ is nonzero for all sufficiently small $t>0$, and
\begin{equation}\label{singlelimit}
 J(e^{2\pi i\alpha},e^{-t/a})
 =e^{2\pi i(a-1)\alpha/a}+O(e^{-c_0/\sqrt t})
\end{equation}
for every $c_0$ with $0<c_0<c_k(\alpha)$, where
\begin{equation}\label{twistgap}
 c_k(\alpha)=\pi^{3/2}\sqrt{\frac{2k}{3}}
       \left(\sqrt{1-|\alpha|}-\sqrt{|\alpha|}\right).
\end{equation}
The estimate is for each fixed $\alpha$ and is not uniform as $|\alpha|\to1/2$.
\end{lemma}
\begin{proof}
Apply Poisson summation to the Schwartz function
$h(x)=e^{2\pi i\alpha x}G(t(x+c))$. Its Fourier transform at an integer $n$ is
$t^{-1}e^{2\pi i(n-\alpha)c}\widehat G((n-\alpha)/t)$, proving \eqref{Poisson}. The hypotheses for this form of Poisson summation are supplied by Lemma~\ref{Gschwartz}; see also \cite{SteinShakarchi}.

The numerator and denominator of $J$ are respectively $S(1/a,\alpha;t)$ and $S(1,\alpha;t)$. If $0<|\alpha|<1/2$, the unique smallest frequency in absolute value is $n=0$. Lemma~\ref{Fourier}, together with $\widehat G(-\xi)=\overline{\widehat G(\xi)}$ for real $\xi$, shows that this term is nonzero for small $t$. The other frequencies satisfy $|n-\alpha|\ge1-|\alpha|>|\alpha|$. Their sum divided by the dominant term is bounded by a power of $t^{-1}$ times
\[
 \exp\!\left[-\frac{\pi^{3/2}\sqrt{2k/3}}{\sqrt t}
             \bigl(\sqrt{1-|\alpha|}-\sqrt{|\alpha|}\bigr)\right].
\]
The remaining sum over $n$ is bounded by comparison with the integral of a polynomial times $e^{-c\sqrt{x/t}}$. Decreasing the positive exponential constant absorbs the power of $t^{-1}$. Division of the dominant phase factors proves \eqref{singlelimit}. If $\alpha=0$, the dominant term is $\widehat G(0)/t>0$, and the same tail estimate applies to $n\ne0$.
\end{proof}

Put
\begin{equation}\label{phase}
 \Theta(t)=\pi^{3/2}\sqrt{\frac{k}{3t}},\qquad
 \phi=\frac{(k-3)\pi}{8}.
\end{equation}
Denote the numerator and denominator of the alternating quotient by
\[
 U_-(t)=qU(-Q,q)=S(1/a,1/2;t),\qquad
 V_-(t)=QV(-Q,q)=S(1,1/2;t).
\]
Then $V_-(t)=Q\mathcal D(Q)$ with $Q=e^{-t}$ and $\mathcal D$ as in \eqref{commonD}; its zero set is independent of $a$ and $d$.

\begin{lemma}\label{alternating}
There is an explicit positive amplitude
\[
 a_k(t)=\frac{2B_k}{t}(2t)^{-(k-3)/4}e^{-\Theta(t)}
\]
such that
\begin{align}
 V_-(t)&=-a_k(t)\{\cos(\Theta(t)-\phi)+\epsilon(t)\},\label{Dminus}\\
 U_-(t)&= a_k(t)\{\cos(\Theta(t)-\phi-\pi/a)+\epsilon_a(t)\},\label{Nminus}
\end{align}
where
\begin{equation}\label{Eest}
 \epsilon(t),\epsilon_a(t)=O(\sqrt t),\qquad \epsilon'(t),\epsilon_a'(t)=O(t^{-1}).
\end{equation}
\end{lemma}
\begin{proof}
In \eqref{Poisson}, the terms $n=0,1$ give
\[
 V_-(t)=-\frac2t\Re\widehat G(1/(2t))+\text{tail},\qquad
 U_-(t)=\frac2t\Re\{e^{i\pi/a}\widehat G(1/(2t))\}+\text{tail}.
\]
The tails involve frequencies $n/(2t)$ with positive odd $n\ge3$. Formula \eqref{Fourierexp} gives \eqref{Dminus}--\eqref{Nminus}. After division by $a_k(t)$, the tails are $O(e^{-c/\sqrt t})$, since their first exponential rate is $\sqrt3\Theta(t)$.

For the derivative estimates, the relative error in the main pair is a real part of $e^{-i(\Theta-\phi)}\varepsilon(1/(2t))$, with a constant phase factor in the numerator. Since $\Theta'(t)=-\Theta(t)/(2t)$, Lemma~\ref{Fourier} bounds its derivative by $O(t^{-1})+O(t^{-1/2})$. Differentiating \eqref{Fourierexp} also gives
\[
 |\widehat G'(\xi)|\le C_k\xi^{(k-3)/4-1/2}
                  e^{-\pi^{3/2}\sqrt{2k\xi/3}}
\]
for large positive $\xi$. This bound justifies termwise differentiation of the tails on compact $t$-intervals and, after division by $a_k$, bounds their derivatives by a power of $t^{-1}$ times $e^{-c/\sqrt t}$. These are $O(t^{-1})$ as $t\downarrow0$.
\end{proof}

\begin{theorem}\label{radial}
Fix integers $a\ge2$, $d\ge2$, and $k\ge1$. If $d$ is odd, then
\begin{equation}\label{oddlimit}
 W_{a,d}(e^{-t/a},k)
 =\frac{\sin(\pi/a)}{d\sin(\pi(a-1)/(ad))}
       +O(e^{-c/\sqrt t}).
\end{equation}
If $d$ is even, then
\begin{equation}\label{tangent}
 W_{a,d}(e^{-t/a},k)
 =\frac{\sin(\pi/a)}d
   \left\{\cot\frac{\pi(a-1)}{ad}-\tan(\Theta(t)-\phi)\right\}
       +O(\sqrt t/\delta^2)
\end{equation}
uniformly for $0<\delta\le1$ on $|\cos(\Theta(t)-\phi)|\ge\delta$, provided $\sqrt t\le c\delta$ for a sufficiently small constant $c=c(a,d,k)>0$. The constant in the $O$ term is independent of $\delta$.
\end{theorem}
\begin{proof}
For odd $d$, use representatives $r=-(d-1)/2,\ldots,(d-1)/2$ in \eqref{Wanalytic} and apply Lemma~\ref{twists}. The sum of limiting phases is
\[
 \sum_r e^{2\pi i(a-1)r/(ad)}
 =\frac{\sin(\pi(a-1)/a)}{\sin(\pi(a-1)/(ad))},
\]
a Dirichlet kernel value. This proves \eqref{oddlimit}.

For even $d$, isolate the alternating term. The sum of the remaining limiting phases, using $r=-d/2+1,\ldots,d/2-1$, is
\[
 \sin(\pi/a)\cot\frac{\pi(a-1)}{ad}+\cos(\pi/a).
\]
By Lemma~\ref{alternating}, away from the indicated phase zeros,
\[
 \frac{U_-}{V_-}
 =-\cos(\pi/a)-\sin(\pi/a)\tan(\Theta-\phi)+O(\sqrt t/\delta^2).
\]
To make the error uniform, write the alternating quotient as $-(u+e_1)/(v+e_2)$, where $u=\cos(\Theta-\phi-\pi/a)$, $v=\cos(\Theta-\phi)$, and $|e_1|+|e_2|\le C\sqrt t$. If $|v|\ge\delta$ and $C\sqrt t\le\delta/2$, subtraction of $-u/v$ bounds the difference by
\[
 \frac{|e_1|}{|v+e_2|}
 +\frac{|u||e_2|}{|v|\,|v+e_2|}
 \le C'\frac{\sqrt t}{\delta^2}.
\]
Adding the nonalternating terms and dividing by $d$ proves \eqref{tangent}.
\end{proof}

\begin{example}
For every positive integer $k$,
\[
 \lim_{q\to1^-}W_{2,3}(q,k)=\frac23,\qquad
 \lim_{q\to1^-}W_{2,5}(q,k)=\frac{1+\sqrt5}{5},\qquad
 \lim_{q\to1^-}W_{4,3}(q,k)=\frac13.
\]
For $a=d=2$, \eqref{tangent} becomes $W_{2,2}=(1-\tan(\Theta-\phi))/2+O(\sqrt t/\delta^2)$.
\end{example}

\begin{theorem}\label{poles}
For fixed $a\ge2$, even $d\ge2$, and $k\ge1$, the function $W_{a,d}(q,k)$ has infinitely many simple real poles approaching $1$. For every sufficiently large integer $j$, there is exactly one such pole $q_j=e^{-t_j/a}$ in the phase window
\[
 |\Theta(t)-\phi-\pi(j+1/2)|<\pi/4.
\]
These are all its real poles sufficiently close to $1$. They satisfy
\begin{align}
 \Theta(t_j)&=\pi\left(j+\frac{k+1}{8}\right)+O(j^{-1}),\label{polephase}\\
 1-q_j&\sim\frac{\pi k}{3aj^2},\label{polelocation}\\
 \Res_{q=q_j}W_{a,d}(q,k)
 &=\frac{2q_jt_j\sin(\pi/a)}{ad\Theta(t_j)}
       (1+O(\sqrt{t_j}))
 \sim\frac{2k\sin(\pi/a)}{3ad\,j^3}.\label{poleresidue}
\end{align}
In particular the residues are eventually positive. In the variable $Q=q^a$, these poles are precisely all sufficiently late zeros of the single function $\mathcal D(Q)$ in \eqref{commonD}. The denominator zeros obey $q_{a,j}=q_{2,j}^{2/a}$ under the common phase indexing.
\end{theorem}
\begin{proof}
Write $\psi=\Theta-\phi$. By \eqref{Dminus}, the zeros of $V_-$ are the zeros of $\cos\psi+\epsilon$. At the two endpoints of each late phase window, $\cos\psi$ has opposite signs and absolute value $1/\sqrt2$, whereas $\epsilon=O(\sqrt t)$. A zero therefore exists. In the window $|\sin\psi|\ge1/\sqrt2$, and
\[
 \frac{d}{dt}(\cos\psi+\epsilon)=-\sin\psi\,\Theta'(t)+O(t^{-1}).
\]
The first term has constant sign and magnitude comparable to $t^{-3/2}$. It dominates the error for small $t$, proving uniqueness and simplicity. Outside the windows, $|\cos\psi|\ge1/\sqrt2$, so there are no other sufficiently small positive zeros.

At a zero, $\cos\psi=O(\sqrt t)$ and $|\sin\psi|=1+O(t)$. Hence \eqref{Nminus} gives
\[
 U_-(t_j)=a_k(t_j)\sin(\pi/a)\sin\psi(t_j)
                  (1+O(\sqrt{t_j}))\ne0.
\]
All the other twisted quotients are holomorphic in a neighborhood of this real point and have nonzero denominators by Lemma~\ref{twists}. Thus the zero is an actual simple pole of the average, with no cancellation.

The zero equation gives $\psi(t_j)=\pi(j+1/2)+O(\sqrt{t_j})$. It first implies $t_j\asymp j^{-2}$, then gives \eqref{polephase} and \eqref{polelocation}. At the zero, the derivative of the amplitude in \eqref{Dminus} contributes nothing, and
\[
 V_-'(t_j)=a_k(t_j)\sin\psi(t_j)\Theta'(t_j)
                         (1+O(\sqrt{t_j})).
\]
It follows that
\[
 \frac{U_-(t_j)}{V_-'(t_j)}
 =-\frac{2t_j\sin(\pi/a)}{\Theta(t_j)}(1+O(\sqrt{t_j})).
\]
Since $dq/dt=-q/a$, conversion to the $q$-residue and the factor $1/d$ in \eqref{Wanalytic} prove \eqref{poleresidue}. The final identity follows because $V_-(t)$ is independent of $a$.
\end{proof}

\begin{corollary}\label{secondpolelaw}
For the zeros indexed in Theorem~\ref{poles}, put
\[
 H=j+\frac{k+1}{8},\qquad
 \alpha_k=\frac{(k-1)(k-3)}{16}
           -\frac{\pi^2k}{6}\left(1-\frac{k}{24}\right).
\]
Then, for fixed $k\ge1$ as $j\to\infty$,
\begin{align}
 \Theta(t_j)&=\pi H-\frac{\alpha_k}{\pi H}+O(H^{-2}),
                                      \label{secondpolephase}\\
 t_j&=\frac{\pi k}{3H^2}
       \left(1+\frac{2\alpha_k}{\pi^2H^2}+O(H^{-3})\right).
                                      \label{secondpolet}
\end{align}
\end{corollary}
\begin{proof}
Retain $A=\pi^2k/6$, $c=1-k/24$, and $b=c+2\pi i\xi$ from the proof of
Lemma~\ref{Fourier}. The uniform expansion
\[
 K_\nu(z)=\sqrt{\frac\pi{2z}}e^{-z}
       \left(1+\frac{4\nu^2-1}{8z}+O(z^{-2})\right)
\]
follows from \cite[\S10.40(i)]{DLMF}. At $\xi=1/(2t)$,
\[
 2\sqrt{Ab}=(1+i)\Theta+\frac{Ac(1-i)}{\Theta}
                         +O(\Theta^{-3}).
\]
The factor $b^{(k-3)/4}$ changes only by $1+O(\Theta^{-2})$ from its
leading term. Substitution in \eqref{I0}, with $\nu=k/2-1$, therefore gives
\[
 \widehat G(1/(2t))=B_k(2t)^{-(k-3)/4}e^{-(1+i)\Theta+i\phi}
       \left(1+\frac{(1-i)\alpha_k}{\Theta}+O(\Theta^{-2})\right).
\]
The product remainder and the other Fourier frequencies are exponentially
smaller, as in Lemmas~\ref{Fourier} and \ref{alternating}. The real-part
zero equation is thus
\[
 \left(1+\frac{\alpha_k}{\Theta}\right)
       \cos\left(\Theta-\phi+\frac{\alpha_k}{\Theta}\right)
       +O(\Theta^{-2})=0.
\]
The unique zero in the $j$th phase window satisfies
$\Theta+\alpha_k/\Theta=\pi H+O(H^{-2})$.
Using $\Theta=\pi H+O(H^{-1})$ proves \eqref{secondpolephase}.
The identity $t=\pi^3k/(3\Theta^2)$ then gives \eqref{secondpolet}.
\end{proof}

For example, $\alpha_3=-7\pi^2/16$, in agreement with
\eqref{refinedphase}, whereas $\alpha_{24}=483/16$.

\begin{corollary}\label{polespacing}
For fixed $a,k$ and even $d$, let $q_j$ be the poles indexed in
Theorem~\ref{poles}. Then
\begin{equation}\label{spacingratio}
 q_{j+1}-q_j\sim\frac{2\pi k}{3aj^3},\qquad
 \frac{\Res_{q=q_j}W_{a,d}(q,k)}{q_{j+1}-q_j}
       \longrightarrow\frac{\sin(\pi/a)}{\pi d}.
\end{equation}
For any fixed $q_0<1$ sufficiently close to $1$,
\begin{equation}\label{polecount}
 \#\{j:q_0\le q_j\le1-\varepsilon\}
 =\sqrt{\frac{\pi k}{3a}}\,\varepsilon^{-1/2}+O(1)
 \qquad(\varepsilon\downarrow0).
\end{equation}
\end{corollary}
\begin{proof}
The phase estimate \eqref{polephase}, the definition of $\Theta$, and
$1-e^{-t/a}=t/a+O(t^2)$ give
\[
 1-q_j=\frac{\pi k}{3a(j+(k+1)/8)^2}+O(j^{-4}).
\]
Subtracting consecutive terms proves the first assertion in
\eqref{spacingratio}; the remainder is smaller than the leading spacing.
The residue asymptotic \eqref{poleresidue} proves the second assertion.
Inverting the same estimate gives
$(1-q_j)^{-1/2}=\sqrt{3a/(\pi k)}\,(j+(k+1)/8)+O(j^{-1})$.
Counting the corresponding integers proves \eqref{polecount}.
\end{proof}

For $a=d=2$, the limiting residue-to-spacing ratio is $1/(2\pi)$,
independently of $k$.

\begin{corollary}\label{nonmodular}
For $a\ge2$, even $d\ge2$, and $k\ge1$, the function
$\tau\mapsto W_{a,d}(e^{2\pi i\tau},k)$ is not a meromorphic modular function on any finite-index subgroup of $\mathrm{SL}_2(\Z)$ with meromorphic behaviour at the cusps.
\end{corollary}
\begin{proof}
The poles in Theorem~\ref{poles} correspond to $\tau_j=it_j/(2\pi a)\to0$. If the function were meromorphic modular, a local parameter at the cusp $0$ would express it as a meromorphic function in a disk about zero; see \cite{DiamondShurman}. Such a function has no sequence of distinct poles tending to the center. The poles $\tau_j$ are distinct in a sufficiently small cusp neighborhood, giving a contradiction.
\end{proof}

\section{The cubic refinement and exact cancellation}
The cubic Euler product admits an exact transform. We include a derivation to fix its normalization and justify the interchange involved. The underlying Jacobi identity is classical; see \cite{AndrewsAskeyRoy,AndrewsBerndt}.

\begin{lemma}\label{cubictransform}
For $k=3$ and real $\xi$,
\begin{equation}\label{exactFourier}
 \widehat G(\xi)=2\pi\sech\!\left(\frac\pi2\sqrt{7+16\pi i\xi}\right).
\end{equation}
Consequently,
\begin{equation}\label{exactD}
 V_-(t)=-\frac{4\pi}{t}
 \sum_{\substack{n\ge1\\n\text{ odd}}}
 \Re\sech\!\left(\frac\pi2\sqrt{7+\frac{8\pi in}{t}}\right),
\end{equation}
with absolute convergence, locally uniform for $t>0$.
\end{lemma}
\begin{proof}
Jacobi's identity is
\[
 \Phi(q)^3=\sum_{n\ge0}(-1)^n(2n+1)q^{n(n+1)/2}.
\]
Replace $q=e^{-v}$ by $e^{-v-\epsilon}$, with $\epsilon>0$, in this identity and integrate against $e^{-v-2\pi i\xi v}$. The integrated series converges absolutely. On the integral side, $0<\Phi(e^{-v-\epsilon})^3\le1$, so dominated convergence applies as $\epsilon\downarrow0$. On the series side, summation by parts shows that the decreasing weights $e^{-\epsilon n(n+1)/2}$ preserve the limit of the convergent alternating series. Hence
\[
 \widehat G(\xi)=2\sum_{n\ge0}
 \frac{(-1)^n(2n+1)}{n^2+n+2+4\pi i\xi}.
\]
Let $\alpha,\beta$ be the roots of $z^2+z+2+4\pi i\xi=0$. The summand without the factor $2$ is invariant under $n\mapsto-1-n$, and
\[
 \frac{2n+1}{(n-\alpha)(n-\beta)}
       =\frac1{n-\alpha}+\frac1{n-\beta}.
\]
The classical partial-fraction expansion
$\sum_{n\in\Z}(-1)^n/(n-z)=-\pi\csc(\pi z)$, interpreted by symmetric limits, now gives
\[
 \widehat G(\xi)=-\pi\{\csc(\pi\alpha)+\csc(\pi\beta)\}.
\]
Since $\beta=-1-\alpha$, the two sine values agree. Taking
$\alpha=(-1+\sqrt{-7-16\pi i\xi})/2$ proves \eqref{exactFourier}. This is also the rescaled eta-cube Laplace transform in \cite[eq.~(14)]{Glasser}; related eta integrals are treated in \cite{Coffey}.

Pairing the positive and negative frequencies in \eqref{Poisson} gives \eqref{exactD}. Its terms have size $O(e^{-c\sqrt{n/t}})$ for large odd $n$, which proves the convergence assertions.
\end{proof}

\begin{theorem}\label{cubicpoles}
Let $k=3$, and index the positive zeros $t_j$ of $V_-$ as in
Theorem~\ref{poles}. Put
\[
 h=j+\frac12,\qquad t_j^*=\frac{\pi}{h\sqrt{h^2+7/4}}.
\]
Then
\begin{equation}\label{cubicpolelaw}
 t_j=t_j^*+O\!\left(h^{-3}e^{-(\sqrt3-1)\pi h}\right),
\end{equation}
and, more precisely,
\begin{equation}\label{cubicsharp}
 (t_j-t_j^*)h^3e^{(\sqrt3-1)\pi h}
 =2(-1)^{j+1}\cos(\sqrt3\pi h)+O(h^{-1}).
\end{equation}
The set of limit points of the left side is $[-2,2]$. In particular,
\[
 \limsup_{j\to\infty}|t_j-t_j^*|h^3e^{(\sqrt3-1)\pi h}=2,
\]
so no bound $O(h^{-3}e^{-ch})$ with $c>(\sqrt3-1)\pi$ can hold for all
sufficiently large $j$. Moreover,
\begin{equation}\label{refinedphase}
 \Theta(t_j)=\pi\left(h+\frac7{16h}+O(h^{-3})\right).
\end{equation}
The corresponding poles of every even section are $q_j=e^{-t_j/a}$.
\end{theorem}
\begin{proof}
Write
\[
 z_n(t)=\frac\pi2\sqrt{7+8\pi in/t},\qquad
 H_n(t)=\Re\sech z_n(t),\qquad n\ge1\text{ odd}.
\]
By \eqref{exactD}, the zero equation is $\sum_{n\ge1\text{ odd}}H_n(t_j)=0$.
If $z_1=u+iv$, then
\begin{equation}\label{sechreal}
 H_1(t)=\frac{\cosh u\cos v}{\sinh^2u+\cos^2v},\qquad
 u^2-v^2=\frac{7\pi^2}{4},\qquad uv=\frac{\pi^3}{t}.
\end{equation}
Consequently
\[
 z_1(t_j^*)=\pi\sqrt{h^2+7/4}+i\pi h,\qquad H_1(t_j^*)=0.
\]
The phase estimate in Theorem~\ref{poles} first gives
$t_j=\pi h^{-2}(1+O(h^{-2}))$, and hence $t_j-t_j^*=O(h^{-4})$.
We may therefore work on a segment $|t-t_j^*|\le Ch^{-4}$ containing
both points. On this segment,
$u=\pi h+O(h^{-1})$ and $v=\pi h+O(h^{-1})$.
Differentiation of \eqref{sechreal} gives
\begin{equation}\label{cubicderivative}
 H_1'(t)=(-1)^j h^3e^{-\pi h}(1+O(h^{-1})),\qquad
 H_1''(t)=O(h^6e^{-\pi h}).
\end{equation}
For the leading constant, at the center one has
\[
 |v'(t_j^*)|^{-1}
 =\frac{2h^2+7/4}{h^2(h^2+7/4)^{3/2}}
 =2h^{-3}(1+O(h^{-2})),
\]
and $\cosh u/\sinh^2u=2e^{-u}(1+O(e^{-2u}))$.

For every odd $n\ge3$, the square-root formula gives, uniformly on the
same segment,
\[
 \Re z_n(t)\ge\sqrt n(\pi h-C/h),\qquad
 |z_n'(t)|\le C\sqrt n\,h^3.
\]
Using $|\sech z|+|\sech z\tanh z|\le Ce^{-\Re z}$ and comparing the
remaining sums with their integrals yields
\begin{equation}\label{cubictail}
 \sum_{n\ge3\text{ odd}}|H_n(t)|=O(e^{-\sqrt3\pi h}),\qquad
 \sum_{n\ge3\text{ odd}}|H_n'(t)|=O(h^3e^{-\sqrt3\pi h}).
\end{equation}
The mean value theorem applied to $H_1$, followed by the zero equation,
now proves \eqref{cubicpolelaw}.

At the next frequency,
$\Re z_3(t_j^*)=\sqrt3\pi h+O(h^{-1})$ and
$\Im z_3(t_j^*)=\sqrt3\pi h+O(h^{-1})$. Therefore
\begin{equation}\label{cubicthird}
 H_3(t_j^*)=2e^{-\sqrt3\pi h}
             \bigl(\cos(\sqrt3\pi h)+O(h^{-1})\bigr).
\end{equation}
The tail in \eqref{cubictail} starting at $n=5$ is
$O(e^{-\sqrt5\pi h})$. Put $\delta=t_j-t_j^*$.
Taylor's theorem and \eqref{cubicderivative}--\eqref{cubictail} give
\[
 0=H_1'(t_j^*)\delta+H_3(t_j^*)+
 O\!\left(e^{-\sqrt5\pi h}
       +h^6e^{-\pi h}\delta^2
       +h^3e^{-\sqrt3\pi h}|\delta|\right).
\]
After division by $H_1'(t_j^*)$, the error is
$O(h^{-4}e^{-(\sqrt3-1)\pi h})$, by \eqref{cubicpolelaw}.
Equations \eqref{cubicderivative} and \eqref{cubicthird} prove
\eqref{cubicsharp}. Since
\[
 (-1)^{j+1}\cos(\sqrt3\pi h)
 =\cos\!\left(\pi(\sqrt3+1)j+\frac{\pi\sqrt3}{2}+\pi\right),
\]
the elementary density of irrational rotations gives the asserted
limit-point set. Finally,
$\Theta(t_j^*)=\pi h(1+7/(4h^2))^{1/4}$;
\eqref{cubicpolelaw} and Taylor expansion prove \eqref{refinedphase}.
\end{proof}

\begin{example}
The elementary part of \eqref{cubicpolelaw} contains all algebraic orders:
\[
 \frac{\pi}{h\sqrt{h^2+7/4}}
 =\frac\pi{h^2}\left(1-\frac7{8h^2}+\frac{147}{128h^4}
             -\frac{1715}{1024h^6}+O(h^{-8})\right).
\]
\end{example}

Primitive roots provide a different consequence of the twisted limit. Recall that $W_{a,1}(q,k)=J(1,q)$. For $d\ge3$, define
\begin{equation}\label{Addef}
 \mathcal A_d(q)=\sum_{e\mid d}\mu(d/e)eW_{a,e}(q,k),
\end{equation}
where $\mu$ is the M\"obius function. The dilation and Euler power remain fixed in this notation.

\begin{corollary}\label{primitive}
For $d\ge3$,
\begin{equation}\label{primitiveidentity}
 \mathcal A_d(q)=\sum_{\ord(z)=d}J(z,q).
\end{equation}
In particular the alternating quotient $J(-1,q)$ cancels identically, and $\mathcal A_d$ has no real poles on some interval ending at $q=1$. Moreover,
\begin{equation}\label{primitivelimit}
 \mathcal A_d(e^{-t/a})=
 2\sum_{\substack{1\le r<d/2\\(r,d)=1}}
 \cos\frac{2\pi(a-1)r}{ad}+O(e^{-c/\sqrt t}).
\end{equation}
\end{corollary}
\begin{proof}
The sum $eW_{a,e}$ contains the contribution of every root whose order divides $e$. In \eqref{Addef}, a root of order $b\mid d$ therefore has coefficient
$\sum_{b\mid e\mid d}\mu(d/e)$, which is $1$ for $b=d$ and $0$ otherwise. This is ordinary M\"obius inversion, as in \cite[Chapter 3]{Stanley}, and proves \eqref{primitiveidentity} as a meromorphic identity. Neither $1$ nor $-1$ is primitive of order $d\ge3$. Lemma~\ref{twists} thus proves that each denominator on the right is nonzero on a sufficiently late real interval. Applying \eqref{singlelimit} and pairing conjugate roots gives \eqref{primitivelimit}.
\end{proof}

\begin{example}\label{primitiveexamples}
For $a=2$ and every fixed integer $k\ge1$, the primitive combinations
in \eqref{Addef} have the following exact limits:
\begin{center}
\begin{tabular}{@{}cl@{\qquad}cl@{}}
\toprule
$d$ & $\displaystyle\lim_{q\to1^-}\mathcal A_d(q)$
& $d$ & $\displaystyle\lim_{q\to1^-}\mathcal A_d(q)$\\
\midrule
3 & $1$ & 12 & $\sqrt6$\\
4 & $\sqrt2$ & 15 & $\sqrt{15+6\sqrt5}$\\
5 & $\sqrt5$ & 16 & $\sqrt{2(2+\sqrt2)(2+\sqrt{2+\sqrt2})}$\\
6 & $\sqrt3$ & 20 & $\dfrac{1+\sqrt5+\sqrt{10+2\sqrt5}}{\sqrt2}$\\[6pt]
8 & $\sqrt{4+2\sqrt2}$ & 24 & $\sqrt{2(2+\sqrt2)(2+\sqrt3)}$\\
10 & $\sqrt{5+2\sqrt5}$ & 30 & $\dfrac{\sqrt3(1+\sqrt5)+\sqrt{10+2\sqrt5}}2$\\[6pt]
\bottomrule
\end{tabular}
\end{center}
The limits use the holomorphic continuations through any removable real
poles of the combinations. In particular, with $(q,k)$ suppressed,
\begin{align*}
 \mathcal A_4&=4W_{2,4}-2W_{2,2},\\
 \mathcal A_{12}&=12W_{2,12}-6W_{2,6}-4W_{2,4}+2W_{2,2},\\
 \mathcal A_{24}&=24W_{2,24}-12W_{2,12}-8W_{2,8}+4W_{2,4}.
\end{align*}
Their limits are independent of the Euler power $k$.

For the evaluations, write
$s_d=2\sum_{1\le r<d/2,\,(r,d)=1}\cos(\pi r/d)$, as in
\eqref{primitivelimit}. The elementary angles and
\[
 2\cos\frac\pi5=\frac{1+\sqrt5}{2},\qquad
 2\cos\frac{2\pi}5=\frac{\sqrt5-1}{2},\qquad
 4\cos\frac\pi{10}=\sqrt{10+2\sqrt5}
\]
give the entries with $d=3,4,5,6$.
The first two identities follow from
$1+2\cos(2\pi/5)+2\cos(4\pi/5)=0$ and the double-angle formula;
the third is the positive half-angle identity.
Sum-to-product and half-angle formulas now give
\begin{align*}
 s_8^2&=4+2\sqrt2,&
 s_{10}^2&=5+2\sqrt5,&
 s_{12}&=\sqrt6,\\
 s_{15}&=\sqrt3\,s_{10},&
 s_{16}&=\frac1{\sin(\pi/16)},\\
 s_{20}&=2\sqrt2\left(\cos\frac\pi5+\cos\frac\pi{10}\right),&&\\
 s_{30}&=2\sqrt3\cos\frac\pi5+2\cos\frac\pi{10}.&&
\end{align*}
For $s_{15}$, pair the indices $1,4$ and $2,7$; for $s_{20}$,
pair $1,9$ and $3,7$; and for $s_{30}$, pair $1,11$ and $7,13$.
The four-term geometric sum gives $s_{16}$, whose square is
$2(2+\sqrt2)(2+\sqrt{2+\sqrt2})$ by rationalizing the half-angle
formula. Finally, pairing $1,11$ and $5,7$ gives
\[
 s_{24}=4\sqrt2\cos\frac\pi8\cos\frac\pi{12}
       =\sqrt{2(2+\sqrt2)(2+\sqrt3)}.
\]
All cosine terms are positive, fixing the signs of the radicals.

More generally, every constant in \eqref{primitivelimit} is a totally
real algebraic integer in $\Q(\zeta_{ad})^+$, where
$\zeta_{ad}=e^{2\pi i/(ad)}$: each paired summand is a root of unity
plus its inverse. These are exact evaluations of the limiting primitive
combinations; no algebraicity of the individual section values is needed.
\end{example}

\begin{corollary}\label{evencancellation}
Fix $a\ge2$ and $k\ge1$, and let $d,e\ge2$ be even. The function
$dW_{a,d}(q,k)-eW_{a,e}(q,k)$ extends holomorphically through every
sufficiently late common real pole. At each such pole,
\begin{equation}\label{exactresidueratio}
 d\,\Res_{q=q_j}W_{a,d}(q,k)
 =e\,\Res_{q=q_j}W_{a,e}(q,k).
\end{equation}
For some $c>0$, its continued values satisfy
\begin{align}\label{evencancellationlimit}
 &dW_{a,d}(e^{-t/a},k)-eW_{a,e}(e^{-t/a},k)\notag\\
 &\qquad=\sin\frac\pi a
   \left\{\cot\frac{\pi(a-1)}{ad}
          -\cot\frac{\pi(a-1)}{ae}\right\}
      +O(e^{-c/\sqrt t})
 \qquad(t\downarrow0).
\end{align}
\end{corollary}
\begin{proof}
In the two root sums \eqref{Wanalytic}, the alternating term $J(-1,q)$
occurs once and cancels exactly in the difference. Every remaining
denominator is nonzero on a sufficiently late real interval by
Lemma~\ref{twists}, so the difference is holomorphic near every pole in
that interval. Taking residues proves \eqref{exactresidueratio}.
For a section of even order $d$, the sum of the remaining limiting
phases is
\[
 \sin(\pi/a)\cot\frac{\pi(a-1)}{ad}+\cos(\pi/a).
\]
Subtract this expression at orders $d$ and $e$ and apply
\eqref{singlelimit} to the finitely many remaining roots. This gives
\eqref{evencancellationlimit}, including at the removable points.
\end{proof}

\section{Approximation to classical modular values}
We now take $F=\Phi^{24}$. The target $T$ is the modular quotient in \eqref{targetmod}. The first theorem gives a family of CM points at which the section convergence has an explicit bound.

\begin{theorem}\label{CMtheorem}
For every integer $a\ge2$, put $q_a=e^{-2\pi/\sqrt a}$ and $\lambda_a=49e^{-4\pi\sqrt a}$. Then $T(q_a)=a^6$, every $W_{a,d}(q_a,24)$ has no pole at $q_a$, and
\begin{equation}\label{CMbound}
 \left|W_{a,d}(q_a,24)-a^6\right|
 \le\frac{4q_a^{2-2a}}{1-q_a}
        \frac{\lambda_a^d}{1-\lambda_a^d}
 \qquad(d\ge2).
\end{equation}
\end{theorem}
\begin{proof}
At $\tau=i/\sqrt a$, one has $a\tau=-1/\tau$. The classical weight-twelve transformation
$\Delta(-1/\tau)=\tau^{12}\Delta(\tau)$ gives
$\Delta(a\tau)=a^{-6}\Delta(\tau)$, hence $T(q_a)=a^6$. This use of the Fricke fixed point is a standard modular evaluation; see \cite{Serre,Zagier}.

For the bound write $q=q_a$, $Q=q^a$, and $r=(49Q)^{-1}$. Then
\[
 0<Q\le e^{-2\pi\sqrt2}<\frac1{7000},\qquad rQ=\frac1{49}.
\]
For $0\le x<1$, the factors in $F(x)=\Phi(x)^{24}$ lie in $[0,1]$. The inequality $1-\prod b_n\le\sum(1-b_n)$ and $1-(1-u)^{24}\le24u$ give
\begin{equation}\label{realproductbound}
 0\le1-F(x)\le\frac{24x}{1-x}.
\end{equation}
Also $F$ is decreasing on $[0,1)$. The cleared denominator has the analytic continuation
\[
 (1-X)V=F(Q)+\sum_{m\ge1}\{F(Q^{m+1})-F(Q^m)\}X^m
\]
for $|X|<Q^{-1}$. On $|X|\le r$, its nonconstant part has absolute value at most
\[
 \frac{24}{1-Q}\sum_{m\ge1}(Qr)^m=\frac1{2(1-Q)}.
\]
Together with \eqref{realproductbound}, this gives
\begin{equation}\label{Vlower}
 |(1-X)V|\ge\frac{1/2-25Q}{1-Q}>0.49.
\end{equation}
For the cleared numerator, the same argument uses
$F(qQ^m)-F(qQ^{m-1})\le24qQ^{m-1}/(1-q)$ and yields
\[
 |(1-X)U|\le1+\frac{24qr}{1-q}\sum_{m\ge0}(Qr)^m
            =1+\frac{q}{2Q(1-q)}.
\]
Thus the quotient $R$ extends holomorphically to $|X|\le r$, and
\begin{equation}\label{Rbound}
 \sup_{|X|=r}|R(X,q)|\le\frac{4q}{Q(1-q)}.
\end{equation}
Here $q/(Q(1-q))\ge1$, so the final numerical constant follows directly from \eqref{Vlower}. The point $X=1$ is removable in the quotient; the cleared expressions define its continuation there.

Cauchy's estimate gives $|C_m(q)|\le4q\,r^{-m}/(Q(1-q))$. Since $Q<r$, the root samples $X=zQ$, $z^d=1$, are inside this disk, and the convergent coefficient filter agrees with \eqref{Wanalytic}. Therefore
\begin{align*}
 |W_{a,d}(q,24)-T(q)|
 &\le q^{1-a}\frac{4q}{Q(1-q)}
                  \sum_{j\ge1}(Q/r)^{jd}\\
 &=\frac{4q^{2-2a}}{1-q}
                  \frac{(49Q^2)^d}{1-(49Q^2)^d},
\end{align*}
which is \eqref{CMbound}.
\end{proof}

The varying-dilation regime has $q_a\to1$ but $q_a^a\to0$, and differs from the fixed-dilation radial limit. Its leading error has an explicit constant and sign.

\begin{proposition}\label{CMleading}
For each fixed integer $d\ge2$, put $q=q_a=e^{-2\pi/\sqrt a}$ and
$Q=q_a^a=e^{-2\pi\sqrt a}$. As the integer $a$ tends to infinity,
\begin{equation}\label{CMleadingerror}
 W_{a,d}(q_a,24)-a^6
 =24q_a(q_a-1)(-23)^{d-2}Q^{2d-2}
       +O_d(a^6Q^{2d-1}).
\end{equation}
Consequently,
\begin{equation}\label{CMsigned}
 W_{a,d}(q_a,24)-a^6
 \sim48\pi(-1)^{d-1}23^{d-2}a^{-1/2}
          e^{-4\pi(d-1)\sqrt a}.
\end{equation}
\end{proposition}
\begin{proof}
Write $F=\Phi^{24}$ and
\[
 (1-X)U=\sum_{m\ge0}A_mX^m,\qquad
 (1-X)V=\sum_{m\ge0}B_mX^m,\qquad
 R=\sum_{m\ge0}C_mX^m.
\]
The Fricke identity in Theorem~\ref{CMtheorem} gives the exact equalities
\[
 F(q)=\frac{a^6Q}{q}F(Q),\qquad C_0=\frac{a^6Q}{q}.
\]
The cleared coefficients are
\begin{align*}
 A_0&=F(q),& A_m&=F(qQ^m)-F(qQ^{m-1}),\\
 B_0&=F(Q),& B_m&=F(Q^{m+1})-F(Q^m)\qquad(m\ge1).
\end{align*}
Since $F(v)=1-24v+O(v^2)$ near zero, these formulas imply
\begin{align*}
 B_0&=1+O(Q),&
 B_m&=24Q^m+O_m(Q^{m+1})\quad(m\ge1),\\
 A_1&=1+O(a^6Q),&
 A_m&=24qQ^{m-1}+O_m(Q^m)\quad(m\ge2).
\end{align*}
All constants here are uniform as $a\to\infty$: the arguments containing
$Q$ lie in a fixed disk about zero, and $q$ is bounded away from zero.
Coefficient comparison gives
\begin{equation}\label{CMrecursion}
 B_0C_m=A_m-\sum_{j=1}^mB_jC_{m-j}.
\end{equation}
It follows that $C_1=1+O(a^6Q)$ and, for each fixed $m\ge2$,
\begin{equation}\label{CMcoefficient}
 C_m=24(q-1)(-23)^{m-2}Q^{m-1}+O_m(a^6Q^m).
\end{equation}
Indeed, in \eqref{CMrecursion} the term $B_{m-1}C_1$ contributes
$24Q^{m-1}$, the term $B_mC_0$ is $O_m(a^6Q^{m+1})$, and the induction
hypothesis for $2\le m-j\le m-1$ gives the leading coefficient
\[
 24(q-1)\left\{1-24\sum_{\ell=0}^{m-3}(-23)^\ell\right\}
 =24(q-1)(-23)^{m-2}.
\]
The sum is empty when $m=2$. All errors have size $O_m(a^6Q^m)$,
including multiplication by $B_0^{-1}=1+O(Q)$, which proves the induction.

The convergent root filter in the proof of Theorem~\ref{CMtheorem} gives
\[
 W_{a,d}(q,24)-a^6=\frac qQ\sum_{j\ge1}C_{jd}Q^{jd}.
\]
The term $j=1$ is \eqref{CMleadingerror}, by \eqref{CMcoefficient}.
To control all remaining terms uniformly, use \eqref{Rbound} on
$|X|=(49Q)^{-1}$. Cauchy's estimate bounds their contribution by
\[
 \frac{4q^2}{Q^2(1-q)}
       \frac{(49Q^2)^{2d}}{1-(49Q^2)^d}
 =O_d\!\left(\sqrt a\,Q^{4d-2}\right)
 =O_d(a^6Q^{2d-1}).
\]
This proves \eqref{CMleadingerror}. Finally,
$q_a(1-q_a)\sim2\pi a^{-1/2}$ and
$a^{13/2}Q\to0$, so the error is smaller than the main term and
\eqref{CMsigned} follows.
\end{proof}

For example, the leading errors for $d=2$ and $d=3$ are respectively
\[
 -48\pi a^{-1/2}e^{-4\pi\sqrt a},\qquad
 1104\pi a^{-1/2}e^{-8\pi\sqrt a}.
\]
These formulas concern increasing dilation. The explicit bound
\eqref{CMbound} remains valid for every $a\ge2$ and $d\ge2$.

For additional CM examples, including negative nomes, the following complex-disk estimate is convenient.

\begin{lemma}\label{smallnome}
If $0<|q|\le1/20$, then all level-two sections with $k=24$ have no pole and
\begin{equation}\label{smallnomebound}
 \left|W_{2,d}(q,24)-q^{-1}\frac{\Phi(q)^{24}}{\Phi(q^2)^{24}}\right|
 \le\frac{16}{|q|}\frac{(3|q|^2)^d}{1-(3|q|^2)^d}
 \qquad(d\ge2).
\end{equation}
\end{lemma}
\begin{proof}
For $|z|<1$, expansion of the product and the triangle inequality give
\begin{equation}\label{complexproductbound}
 |\Phi(z)^{24}|\le e^{24|z|/(1-|z|)},\qquad
 |\Phi(z)^{24}-1|\le e^{24|z|/(1-|z|)}-1.
\end{equation}
Put $u=|q|$. On $|X|\le1/3$, the original denominator in \eqref{Rdef} satisfies
\[
 |V(X,q)|\ge2-e^{24u^2/(1-u^2)}
              -\frac12e^{24u^4/(1-u^4)}>0.43.
\]
The last inequality follows by using $u\le1/20$; for example
$e^{24/399}<1.063$ and $e^{24/159999}<1.001$ suffice. The numerator satisfies
\[
 |U(X,q)|\le\frac32e^{24u/(1-u)}<6.
\]
Consequently $|R(X,q)|<16$ on this disk. Cauchy's estimate gives $|C_m(q)|\le16\cdot3^m$. The samples $X=zq^2$ lie inside the disk, so \eqref{rootfilter} applies analytically. Summing the selected coefficient bound proves \eqref{smallnomebound}.
\end{proof}

Let $u(\tau)=\Delta(\tau)/\Delta(2\tau)$. The classical level-two relation is
\begin{equation}\label{jrelation}
 j(\tau)=\frac{(u(\tau)+256)^3}{u(\tau)^2}.
\end{equation}
It is Maier's rational parametrization \cite[eq.~(1.6)]{Maier} after replacing his parameter $t_2$ by $4096/u$. The transformation of $\Delta$ and the invariance of $j$ give, as in Corollary~\ref{modularthreshold},
\begin{equation}\label{CMdoubling}
 u\!\left(-\frac1{2\tau}\right)=\frac{4096}{u(\tau)},
 \qquad j(2\tau)=\frac{(u(\tau)+16)^3}{u(\tau)}.
\end{equation}
These classical identities give the exact targets below; the individual section values are not asserted to be algebraic.

\begin{proposition}\label{CMtable}
For every point in the table, with $q=e^{2\pi i\tau}$,
$W_{2,d}(q,24)\to u(\tau)$ as $d\to\infty$, with the bound \eqref{smallnomebound}.
\begin{center}
\begin{tabular}{@{}ll@{\qquad}ll@{}}
\toprule
$\tau$ & $u(\tau)$ & $\tau$ & $u(\tau)$\\
\midrule
$i/2$ & $8$ & $2i$ & $1024\sqrt2(1+\sqrt2)^6$\\
$i/\sqrt2$ & $64$ & $i\sqrt7$ & $4096(8+3\sqrt7)^3$\\
$i\sqrt3/2$ & $4(2+\sqrt3)^3$ & $(1+i)/2$ & $-64$\\
$i$ & $512$ & $(1+i\sqrt2)/2$ & $-8(1+\sqrt2)^3$\\
$i\sqrt7/2$ & $(8+3\sqrt7)^3$ & $(1+i\sqrt3)/2$ & $-256$\\
$i\sqrt2$ & $512(1+\sqrt2)^3$ & $(1+i\sqrt7)/2$ & $-4096$\\
$i\sqrt3$ & $1024(2+\sqrt3)^3$ & &\\
\bottomrule
\end{tabular}
\end{center}
The same conclusion holds at the following three points:
\begin{align*}
 u\left(\tfrac12+i\right)
   &=-2\sqrt2(1+\sqrt2)^6,\\
 u\left(\tfrac12+i\sqrt3\right)
   &=-(3\sqrt6-5\sqrt2)
       \bigl[(1+\sqrt2)(\sqrt2+\sqrt3)\bigr]^6,\\
 u(2i\sqrt3)
   &=1024(3\sqrt6+5\sqrt2)
       \bigl[(1+\sqrt2)(\sqrt2+\sqrt3)\bigr]^6.
\end{align*}
\end{proposition}
\begin{proof}
The classical CM values needed here are
\begin{gather*}
 j(i)=1728,\quad j(i\sqrt2)=8000,\quad j(i\sqrt3)=54000,\\
 j((1+i\sqrt3)/2)=0,\qquad j((1+i\sqrt7)/2)=-3375;
\end{gather*}
see \cite[\S6.1, p.~71]{Zagier} and \cite{Cox}. Substitution into \eqref{jrelation} gives
\begin{align*}
 (u+256)^3-1728u^2&=(u-512)^2(u+64),\\
 (u+256)^3-8000u^2&=(u-64)(u^2-7168u-262144),\\
 (u+256)^3-54000u^2&=(u+16)(u^2-53248u+1048576),\\
 (u+256)^3+3375u^2&=(u+4096)(u^2+47u+4096).
\end{align*}
For real $-1<q<1$, $q\ne0$, the product
\[
 u(\tau)=q^{-1}\prod_{n\ge1}(1+q^n)^{-24}
\]
has the sign of $q$. When $0<q<1$, it also gives
\begin{equation}\label{CMrootbound}
 u(\tau)\ge q^{-1}\exp\!\left(-\frac{24q}{1-q}\right).
\end{equation}
Thus $u(i)=512$. At $i\sqrt2$ and $i\sqrt3$, the nomes are less than $1/1000$ and $1/10000$, respectively, so \eqref{CMrootbound} selects the large positive quadratic roots. When $j=0$, necessarily $u=-256$. At $(1+i\sqrt7)/2$, the quadratic factor has negative discriminant, so reality gives $u=-4096$. The Fricke identity in \eqref{CMdoubling} now gives $u(i/2)=8$ and, at its fixed point, $u(i/\sqrt2)=64$.

For $\tau=(1+i)/2$, its Fricke image is $\tau-1$. Periodicity and \eqref{CMdoubling} give $u(\tau)^2=4096$; the negative nome selects $u=-64$. Moreover, applying the second identity in \eqref{CMdoubling} at $i$ and $(1+i\sqrt7)/2$ gives
\[
 j(2i)=287496,\qquad j(i\sqrt7)=16581375.
\]
The remaining values follow from the exact factorizations
\begin{align*}
 (u+16)^3-8000u&=(u-64)(u^2+112u-64),\\
 (u+16)^3-54000u&=(u+256)(u^2-208u+16),\\
 (u+16)^3-16581375u&=(u+4096)(u^2-4048u+1),\\
 (u+256)^3-287496u^2&=(u-8)(u^2-286720u-2097152),\\
 (u+256)^3-16581375u^2&=(u+1)(u^2-16580608u+16777216).
\end{align*}
The first identity has exactly one negative root, $-56-40\sqrt2$, which selects $\tau=(1+i\sqrt2)/2$. For the next two identities, the positive nomes at $i\sqrt3/2$ and $i\sqrt7/2$ are less than $1/200$; hence \eqref{CMrootbound} gives $u>100$ and selects the large positive root. At $2i$ and $i\sqrt7$, the nomes are less than $1/100000$, so the same bound gives $u>90000$ and again selects the large positive root. Expanding the powers in the table gives precisely these roots.

Splitting the product into even and odd factors gives
\begin{equation}\label{CMhalfperiod}
 u(\tau)u\left(\tau+\tfrac12\right)=-u(2\tau).
\end{equation}
Indeed, with $q=e^{2\pi i\tau}$,
$\prod_{n\text{ odd}}(1-q^{2n})
 =\prod_{n\ge1}(1+q^{2n})^{-1}$, and the identity follows from
$u(\tau)=q^{-1}\prod_{n\ge1}(1+q^n)^{-24}$.
The values at $i$ and $2i$ therefore give the first additional value.

For $y>0$, put $v=u(iy)>0$ and $w=u(\tfrac12+iy)<0$.
Since $j(2iy+1)=j(2iy)$, equation~\eqref{CMdoubling} gives
\[
 (w+16)^3-\frac{(v+16)^3}{v}w
 =(w-v)\left(w^2+(v+48)w-\frac{4096}{v}\right)=0.
\]
The quadratic factor has exactly one negative root. When $y=\sqrt3$,
$v=1024(2+\sqrt3)^3$, and that root is
\[
 w=-13336-9405\sqrt2-7680\sqrt3-5445\sqrt6.
\]
Here $v+48=2(13336+7680\sqrt3)$, and the identity
\[
 (9405\sqrt2+5445\sqrt6)^2-(13336+7680\sqrt3)^2
 =104-60\sqrt3=\frac{4096}{v}
\]
verifies the quadratic equation. Expanding the second additional expression gives this value of $w$.
Finally, \eqref{CMhalfperiod} and
$(2+\sqrt3)^3(3\sqrt6-5\sqrt2)=3\sqrt6+5\sqrt2$
give the value at $2i\sqrt3$.

Every point displayed in the proposition has imaginary part at least
$1/2$, so $0<|q|\le e^{-\pi}<1/20$.
Lemma~\ref{smallnome} proves all the stated limits with the same explicit
error bound.
\end{proof}

More general CM targets can be specified without radical expressions. If $D<0$ is a quadratic discriminant and $H_D$ is the class polynomial of degree $h$, then every corresponding value of $u$ satisfies
\begin{equation}\label{classpoly}
 u^{2h}H_D\!\left(\frac{(u+256)^3}{u^2}\right)=0.
\end{equation}
Indeed $H_D(j(\tau))=0$ by the classical theory \cite{Cox}, and \eqref{jrelation} gives the assertion. The left side is a monic polynomial in $\Z[u]$ of degree $3h$: its leading term comes from $(u+256)^{3h}$, and all remaining terms have smaller degree. It need not be irreducible. Class polynomials may be computed by complex approximation or Chinese remainder methods \cite{Enge,Sutherland}; exact factorization and root isolation then specify the chosen target. Lemma~\ref{smallnome} applies whenever the selected nome lies in its disk.

\section*{Conclusion}
The root-of-unity sections have exact formal contact with their product
quotient, but their behavior near $q=1$ depends on the section order.
The real-pole theorem, the sharp cubic correction, and the signed CM
asymptotic quantify these different regimes.

The main remaining question is global. Does the local sampling pole remain
simple and real until it reaches $z=-1$, and can every competing pole be
excluded before that crossing? An exact convergence threshold also requires
numerator noncancellation. For odd sections, the zero-free results near the
two real endpoints leave the intervening interval unresolved.

Other radial directions require a separate analysis of the product phases
by residue classes. The positive-radius results establish neither a quantum
modular transformation law nor a natural boundary on the whole unit circle.
For a finite product $\prod_\delta\Phi(q^\delta)^{r_\delta}$,
Lemma~\ref{Gschwartz} shows that the exponential behavior is governed by
the weighted sum $\sum_\delta r_\delta/\delta$.
Positive total exponent alone is insufficient. For example,
\[
 e^{-v}\frac{\Phi(e^{-2v})^3}{\Phi(e^{-v})^2}
 \sim\frac{\sqrt\pi}{2\sqrt v}\exp\!\left(\frac{\pi^2}{12v}\right),
 \qquad v\downarrow0,
\]
since the exponents $r_2=3$, $r_1=-2$ give weighted sum $-1/2$.

\section*{Acknowledgment}
The author thanks Prof. George E. Andrews for his encouraging and
supportive words and for the inspiration he has provided.

\begin{thebibliography}{99}
\bibitem{Andrews}
G. E. Andrews, \emph{The Theory of Partitions}, Cambridge University Press, Cambridge, 1984, reprint of the 1976 original.
\href{https://doi.org/10.1017/CBO9780511608650}{doi:10.1017/CBO9780511608650}.

\bibitem{AndrewsAskeyRoy}
G. E. Andrews, R. Askey, and R. Roy, \emph{Special Functions}, Encyclopedia of Mathematics and its Applications 71, Cambridge University Press, Cambridge, 1999.
\href{https://doi.org/10.1017/CBO9781107325937}{doi:10.1017/CBO9781107325937}.

\bibitem{AndrewsBerndt}
G. E. Andrews and B. C. Berndt, \emph{Ramanujan's Lost Notebook, Part II}, Springer, New York, 2009.
\href{https://doi.org/10.1007/b13290}{doi:10.1007/b13290}.

\bibitem{Apostol}
T. M. Apostol, \emph{Modular Functions and Dirichlet Series in Number Theory}, second ed., Graduate Texts in Mathematics 41, Springer, New York, 1990.
\href{https://doi.org/10.1007/978-1-4612-0999-7}{doi:10.1007/978-1-4612-0999-7}.

\bibitem{Coffey}
M. W. Coffey, \emph{Some integrals of the Dedekind $\eta$ function}, 2019.
\href{https://arxiv.org/abs/1901.07168}{arXiv:1901.07168}.

\bibitem{Comtet}
L. Comtet, \emph{Advanced Combinatorics: The Art of Finite and Infinite Expansions}, revised and enlarged ed., D. Reidel, Dordrecht, 1974.
\href{https://doi.org/10.1007/978-94-010-2196-8}{doi:10.1007/978-94-010-2196-8}.

\bibitem{Cox}
D. A. Cox, \emph{Primes of the Form $x^2+ny^2$: Fermat, Class Field Theory, and Complex Multiplication}, second ed., Wiley, Hoboken, NJ, 2013.
\href{https://doi.org/10.1002/9781118400722}{doi:10.1002/9781118400722}.

\bibitem{DiamondShurman}
F. Diamond and J. Shurman, \emph{A First Course in Modular Forms}, Graduate Texts in Mathematics 228, Springer, New York, 2005.
\href{https://doi.org/10.1007/978-0-387-27226-9}{doi:10.1007/978-0-387-27226-9}.

\bibitem{Enge}
A. Enge, The complexity of class polynomial computation via floating point approximations, \emph{Math. Comp.} \textbf{78} (2009), no.~266, 1089--1107.
\href{https://arxiv.org/abs/cs/0601104}{arXiv:cs/0601104}.

\bibitem{FlajoletSedgewick}
P. Flajolet and R. Sedgewick, \emph{Analytic Combinatorics}, Cambridge University Press, Cambridge, 2009.
\href{https://doi.org/10.1017/CBO9780511801655}{doi:10.1017/CBO9780511801655}.

\bibitem{GasperRahman}
G. Gasper and M. Rahman, \emph{Basic Hypergeometric Series}, second ed., Encyclopedia of Mathematics and its Applications 96, Cambridge University Press, Cambridge, 2004.
\href{https://doi.org/10.1017/CBO9780511526251}{doi:10.1017/CBO9780511526251}.

\bibitem{Glasser}
M. L. Glasser, Some integrals of the Dedekind eta-function, \emph{J. Math. Anal. Appl.} \textbf{354} (2009), no.~2, 490--493.
\href{https://doi.org/10.1016/j.jmaa.2009.01.024}{doi:10.1016/j.jmaa.2009.01.024};
\href{https://arxiv.org/abs/0812.1992}{arXiv:0812.1992}.

\bibitem{HardyRamanujan}
G. H. Hardy and S. Ramanujan, Asymptotic formulae in combinatory analysis, \emph{Proc. London Math. Soc.} (2) \textbf{17} (1918), 75--115.
\href{https://doi.org/10.1112/plms/s2-17.1.75}{doi:10.1112/plms/s2-17.1.75}.

\bibitem{KubertLang}
D. S. Kubert and S. Lang, \emph{Modular Units}, Grundlehren der mathematischen Wissenschaften 244, Springer, New York, 1981.
\href{https://doi.org/10.1007/978-1-4757-1741-9}{doi:10.1007/978-1-4757-1741-9}.

\bibitem{Maier}
R. S. Maier, On rationally parametrized modular equations, \emph{J. Ramanujan Math. Soc.} \textbf{24} (2009), 1--73.
\href{https://arxiv.org/abs/math/0611041}{arXiv:math/0611041}.

\bibitem{Melman}
A. Melman, A unifying framework for generalizations of the Enestr\"om--Kakeya theorem, \emph{Elem. Math.} \textbf{75} (2020), no.~1, 1--14.
\href{https://doi.org/10.4171/EM/398}{doi:10.4171/EM/398};
\href{https://arxiv.org/abs/1711.09517}{arXiv:1711.09517}.

\bibitem{DLMF}
NIST, \emph{Digital Library of Mathematical Functions}, Sections 10.32 and 10.40, integral representations and large-argument expansions of modified Bessel functions.
\url{https://dlmf.nist.gov/10.32}; \url{https://dlmf.nist.gov/10.40}.

\bibitem{Olver}
F. W. J. Olver, \emph{Asymptotics and Special Functions}, A K Peters, Wellesley, MA, 1997, reprint of the 1974 original.
\href{https://www.routledge.com/Asymptotics-and-Special-Functions/Olver/p/book/9781568810690}{ISBN 9781568810690}.

\bibitem{Rademacher}
H. Rademacher, \emph{Topics in Analytic Number Theory}, Grundlehren der mathematischen Wissenschaften 169, Springer, Berlin, 1973.
\href{https://doi.org/10.1007/978-3-642-80615-5}{doi:10.1007/978-3-642-80615-5}.

\bibitem{RaghavaEarlier}
K. Srinivasa Raghava, \emph{Evenized dilation quotients from Euler products: a modular skeleton at $k=24N$ and certified dominant pole pairs at $k=24,48,96$}, author-posted preprint record, January 2026.
\href{https://www.researchgate.net/publication/400079147}{ResearchGate record 400079147}.

\bibitem{Schmidt}
M. D. Schmidt, A catalog of interesting and useful Lambert series identities, \emph{Online J. Anal. Comb.} \textbf{21} (2026), Paper 1, 1--22.
\href{https://doi.org/10.61091/ojac21-01}{doi:10.61091/ojac21-01};
\href{https://arxiv.org/abs/2004.02976v2}{arXiv:2004.02976v2}.

\bibitem{Serre}
J.-P. Serre, \emph{A Course in Arithmetic}, Graduate Texts in Mathematics 7, Springer, New York, 1973.
\href{https://doi.org/10.1007/978-1-4684-9884-4}{doi:10.1007/978-1-4684-9884-4}.

\bibitem{Stanley}
R. P. Stanley, \emph{Enumerative Combinatorics, Volume 1}, second ed., Cambridge Studies in Advanced Mathematics 49, Cambridge University Press, Cambridge, 2011.
\href{https://doi.org/10.1017/CBO9781139058520}{doi:10.1017/CBO9781139058520}.

\bibitem{SteinShakarchi}
E. M. Stein and R. Shakarchi, \emph{Fourier Analysis: An Introduction}, Princeton Lectures in Analysis 1, Princeton University Press, Princeton, NJ, 2003.
\href{https://press.princeton.edu/books/hardcover/9780691113845/fourier-analysis}{ISBN 9780691113845}.

\bibitem{Sutherland}
A. V. Sutherland, Computing Hilbert class polynomials with the Chinese remainder theorem, \emph{Math. Comp.} \textbf{80} (2011), no.~273, 501--538.
\href{https://doi.org/10.1090/S0025-5718-2010-02373-7}{doi:10.1090/S0025-5718-2010-02373-7};
\href{https://arxiv.org/abs/0903.2785}{arXiv:0903.2785}.

\bibitem{TrefethenWeideman}
L. N. Trefethen and J. A. C. Weideman, The exponentially convergent trapezoidal rule, \emph{SIAM Rev.} \textbf{56} (2014), no.~3, 385--458.
\href{https://doi.org/10.1137/130932132}{doi:10.1137/130932132}.

\bibitem{Wong}
R. Wong, \emph{Asymptotic Approximations of Integrals}, Classics in Applied Mathematics 34, Society for Industrial and Applied Mathematics, Philadelphia, PA, 2001.
\href{https://doi.org/10.1137/1.9780898719260}{doi:10.1137/1.9780898719260}.

\bibitem{Zagier}
D. Zagier, Elliptic modular forms and their applications, in \emph{The 1-2-3 of Modular Forms}, Universitext, Springer, Berlin, 2008, pp.~1--103.
\href{https://doi.org/10.1007/978-3-540-74119-0_1}{doi:10.1007/978-3-540-74119-0\_1}.

\end{thebibliography}
\end{document}
